\chapter{Traces, determinants and the index}\label{ch:determinants}
\providecommand{\Gs}[1]{\Gamma^{#1}}
\providecommand{\Psiisoph}[1]{\Psi^{#1}_{\mathrm{iso,ph}}}
\providecommand{\detFr}{\operatorname{det}_{\mathrm{Fr}}}
\providecommand{\TrRes}{\operatorname{Tr}_{\mathrm{R}}}
\providecommand{\TrReg}[1]{\overline{\operatorname{Tr}}_{#1}}
\providecommand{\Giso}[1]{G^{#1}_{\mathrm{iso}}}
\providecommand{\FP}{\operatorname*{FP}}
\providecommand{\Ell}{\operatorname{Ell}}

For a linear map between finite-dimensional spaces three numbers are
always available: the trace, the determinant, and the index
$\dim(\text{domain})-\dim(\text{target})$. The trace vanishes on
commutators, the determinant is multiplicative and detects invertibility,
and the index of a square matrix is zero. This chapter asks how much of this
survives for operators on $L^2(\R^n)$, and finds that the failures are as
interesting as the successes: the trace does not extend to all operators,
and the precise way in which it fails computes the index.

In this chapter we use the isotropic calculus of \cref{ch:isotropic} as a
laboratory for three closely related constructions: traces, determinants and
the index. On the residual algebra $\Psiiso{-\infty}(\R^n)$, whose elements
have Schwartz kernels, there is a unique trace, and it gives rise to the
Fredholm determinant on $\Id+\Psiiso{-\infty}(\R^n)$; the determinant leads
directly to analytic Fredholm theory and so to the meromorphy of the
resolvent of an elliptic operator. On the full algebra of polyhomogeneous
operators of integral order the trace no longer extends, but complex powers
allow us to define the \emph{residue trace}, a genuine trace which only sees
one term of the symbol expansion, and the \emph{regularized trace}, which
extends the trace but fails to be a trace; the failure is measured by the
residue trace and an exterior derivation. These ingredients give a formula for
the index of an elliptic operator, an obstruction to inverting families of
elliptic operators by smoothing perturbations, and finally the determinant
line bundle of such a family, together with a natural connection whose
curvature is that obstruction. \Cref{determinants:fig:roadmap} shows how
these constructions depend on each other.

\begin{figure}[htbp]
\centering
\begin{tikzcd}[column sep=1.6em,row sep=1.6em,cells={font=\small}]
  \text{trace }\Tr\text{ on }\Psiiso{-\infty}\arrow[r]\arrow[d]
  & \text{Fredholm determinant }\detFr\arrow[r]
  & \text{analytic Fredholm}\arrow[d]\\
  \text{complex powers }H(z)\arrow[d]\arrow[r]
  & \text{residue trace }\TrRes\arrow[d]
  & \text{resolvent and spectrum}\\
  \text{regularized trace }\TrReg{H}\arrow[r]
  & \text{trace defect formula}\arrow[r]\arrow[d]
  & \text{index formula}\\
  & \text{variation form }\beta\arrow[r]
  & \text{determinant bundle}
\end{tikzcd}
\caption[Roadmap of the chapter on traces and determinants]{Roadmap of the
chapter. The top row uses only the residual algebra; the lower rows use
complex powers to go beyond it.}
\label{determinants:fig:roadmap}
\end{figure}

\section{Setup}\label{determinants:sec:setup}

We recall from \cref{ch:isotropic} what we need, and fix notation. Write
$Z=(x,\xi)\in\R^{2n}$. For $m\in\R$ the space of \emph{isotropic symbols} of
order $m$ is
\begin{equation*}
  \Gs{m}(\R^{2n})=\bigl\{a\in\CI(\R^{2n});\ \abs{D_Z^\alpha a(Z)}\le
  C_\alpha\ang{Z}^{m-\abs{\alpha}}\ \forall\,\alpha\bigr\},
\end{equation*}
a Fr\'echet space with the best constants $C_\alpha$ as seminorms, and
$\Psiiso{m}(\R^n)=\Op(\Gs{m}(\R^{2n}))$, with the left quantization of
the Introduction. The polyhomogeneous subspace $\Psiisoph{m}(\R^n)$ consists of
the operators whose left symbol has an expansion $a\sim\sum_{j\ge0}a_{m-j}$
with $a_{m-j}$ homogeneous of degree $m-j$ in $\abs{Z}\ge1$; equivalently,
$\ang{Z}^{-m}a$ is smooth on the radial compactification $\RC{\R^{2n}}$, i.e.\
in polar coordinates $Z=\theta/t$, $t=\abs{Z}^{-1}$, it is smooth up to $t=0$.
The same definition with $m$ replaced by $z\in\C$ gives the operators
$\Psiisoph{z}(\R^n)\subset\Psiiso{\Re z}(\R^n)$ of complex order. We write
$\Psiisoph{\Z}(\R^n)=\bigcup_{m\in\Z}\Psiisoph{m}(\R^n)$. The \emph{residual}
operators are
\begin{equation*}
  \Psiiso{-\infty}(\R^n)=\bigcap_m\Psiiso{m}(\R^n)=\bigcap_m\Psiisoph{m}(\R^n),
\end{equation*}
which are exactly the operators with Schwartz kernels in $\cS(\R^{2n})$; we
identify $\Psiiso{-\infty}(\R^n)$ with $\cS(\R^{2n})$, including its
topology.

All of this applies to operators on $\C^N$-valued functions:
$\Psiisoph{m}(\R^n;\C^N)$ consists of $N\times N$ matrices of elements of
$\Psiisoph{m}(\R^n)$ (see \cref{determinants:ex:bundles}). Everything below
holds in this generality, with $\tr$ denoting the trace of matrices; to keep
the notation light we often suppress $\C^N$.

We use the following results from \cref{ch:isotropic}.
\begin{enumerate}[label=(I\arabic*)]
\item\label{determinants:I:algebra} \emph{Algebra.} Composition maps
  $\Psiiso{m}\times\Psiiso{m'}$ to $\Psiiso{m+m'}$ continuously for
  the symbol topologies, and $\Psiisoph{z}\times\Psiisoph{z'}$ to
  $\Psiisoph{z+z'}$; the left symbol of a composite is given by
  $\sigma_L(AB)\sim\sum_\alpha\frac1{\alpha!}\pa_\xi^\alpha a\,D_x^\alpha b$.
  Adjoints and transposes of elements of $\Psiiso{m}$ (resp.\ $\Psiisoph{z}$)
  lie in the same space. In particular $\Psiiso{-\infty}$ is a two-sided ideal.
  The principal symbol is multiplicative.
\item\label{determinants:I:bounded} \emph{Boundedness.} Every element of
  $\Psiiso{m}$ is continuous on $\cS(\R^n)$ and on $\cSp(\R^n)$; elements
  of $\Psiiso{0}$ are bounded on $L^2(\R^n)$, with norm bounded by a symbol
  seminorm; elements of negative order are compact on $L^2$, and residual
  operators map $\cSp(\R^n)$ continuously into $\cS(\R^n)$.
\item\label{determinants:I:elliptic} \emph{Ellipticity.} An element
  $A\in\Psiisoph{z}(\R^n;\C^N)$ is elliptic if its principal symbol is invertible
  on the sphere. It then has a parametrix $B\in\Psiisoph{-z}$, with $AB-\Id$ and
  $BA-\Id$ residual; any two parametrices differ by a residual operator.
  If $Au\in\cS$ for some $u\in\cSp$, then $u\in\cS$. There is a
  \emph{generalized inverse} $B\in\Psiisoph{-z}$ with
  \begin{equation}\label{determinants:eq:gen-inverse}
    AB=\Id-\Pi_1,\quad BA=\Id-\Pi_0,
  \end{equation}
  where $\Pi_0$, $\Pi_1$ are the $L^2$-orthogonal projections onto the finite
  dimensional spaces $\Nul A\subset\cS$ and $\Nul A^*\subset\cS$ (null spaces
  on $\cSp$). Replacing $B$ by $(\Id-\Pi_0)B(\Id-\Pi_1)$, which changes
  neither identity since $A\Pi_0=0$ and $\Pi_1A=0$, we may and do assume
  $\Pi_0B=0=B\Pi_1$. In particular an elliptic $A$ which is a bijection on
  $\cSp$ (we then say $A$ is \emph{invertible}) has inverse in $\Psiisoph{-z}$.
  For order $0$, $A$ is invertible in this sense if and only if it is
  invertible on $L^2$.
\item\label{determinants:I:residual} \emph{The residual algebra.} If
  $E\in\Psiiso{-\infty}(\R^n)$ then $\Id+E$ is Fredholm on $L^2$ of index
  zero; if it is invertible on $L^2$ its inverse is $\Id+F$ with
  $F\in\Psiiso{-\infty}$. If $\Nul(\Id+E)\ne\{0\}$ there is a finite rank
  $F'\in\Psiiso{-\infty}$, mapping $\Nul(\Id+E)$ isomorphically onto
  $\Nul(\Id+E^*)$ and vanishing on its orthocomplement, such that $\Id+E+sF'$
  is invertible for all $s\ne0$. (This is also the special case
  $A=\Id+E$ of \cref{determinants:prop:invertible-perturbation} below.) The
  series $\exp E=\sum_k E^k/k!$ converges in $\Id+\Psiiso{-\infty}$.
\item\label{determinants:I:oscillator} \emph{The harmonic oscillator.} The
  eigenfunctions $u_\alpha=c_\alpha C^\alpha e^{-\abs{x}^2/2}$,
  $\alpha\in\N_0^n$, of $H=\abs{D}^2+\abs{x}^2$ (with $C_j=D_j+ix_j$ and
  normalizing constants $c_\alpha>0$) are real-valued and form an orthonormal
  basis of $L^2(\R^n)$; $Hu_\alpha=(n+2\abs{\alpha})u_\alpha$; and each Schwartz
  seminorm of $u_\alpha$ is bounded by a polynomial in $\abs{\alpha}$. See
  also \cref{determinants:ex:oscillator}.
\end{enumerate}

\begin{definition}\label{determinants:def:families}
  Let $Y$ be a manifold and $\Omega\subset\C$ open.
  \begin{enumerate}[label=(\alph*)]
  \item A family $E_y\in\Psiiso{-\infty}(\R^n)$, $y\in Y$, is \emph{smooth}
    if its kernel $E(y,x,x')$ is in $\CI(Y\times\R^{2n})$ and, in local
    coordinates, $\ang{(x,x')}^k\pa_y^\gamma\pa_{x,x'}^\alpha E$ is bounded on
    $K\times\R^{2n}$ for every compact $K$ and all $k,\gamma,\alpha$. A family
    $E(z)$, $z\in\Omega$, is \emph{holomorphic} if it is smooth and $E$ is
    holomorphic in $z$.
  \item Let $h$ be holomorphic on $\Omega$. A family $A(z)=\Op(a(z))$,
    $z\in\Omega$, is a \emph{holomorphic family of order $h(z)$} if $a\in\CI(\Omega\times\R^{2n})$ is holomorphic in $z$ and
    \begin{equation}\label{determinants:eq:hol-family}
      a(z;r\theta)=r^{h(z)}\,b(z;1/r,\theta),\quad r\ge1,\ \theta\in\Sphere^{2n-1},
    \end{equation}
    with $b\in\CI(\Omega\times[0,1]\times\Sphere^{2n-1})$ holomorphic in $z$.
    Then $A(z)\in\Psiisoph{h(z)}$, and the term of homogeneity $h(z)-j$ in the
    expansion of $a(z)$ is $\abs{Z}^{h(z)-j}b_j(z;Z/\abs{Z})$, with
    $b_j(z;\theta)=\frac1{j!}\pa_t^jb(z;0,\theta)$. Families of fixed order
    $m$ depending smoothly on $y\in Y$ are defined in the same way, with $b$
    smooth in $(y,t,\theta)$.
  \end{enumerate}
\end{definition}

\begin{enumerate}[resume*]
\item\label{determinants:I:families} \emph{Families.} Composites of
  holomorphic (or smooth) families are holomorphic (or smooth), with orders
  adding; this includes composites with residual families. If $A(z)$ is a
  holomorphic family of elliptic operators of order $h(z)$, then near each
  point of $\Omega$ there is a holomorphic family $B(z)$ of order $-h(z)$ such
  that $A(z)B(z)-\Id$ and $B(z)A(z)-\Id$ are holomorphic families in
  $\Psiiso{-\infty}$ (the parametrix construction of \cref{ch:isotropic}, with
  Borel's lemma applied in the variable $t$ uniformly in $z$); the same holds
  for smooth families. Finally
  \begin{equation}\label{determinants:eq:standard-regularizer}
    H(z)=\Op(\ang{Z}^z)\,\Id_N
  \end{equation}
  is an entire holomorphic family of order $z$ with $H(0)=\Id$.
\end{enumerate}

We complete these facts with three elementary lemmas about residual
operators. Fix seminorms on $\cS(\R^{2n})$ by
\begin{equation*}
  \norm{K}_{(k)}=\sum_{\abs{\alpha}+\abs{\beta}+\abs{\gamma}+\abs{\delta}\le k}
  \sup_{x,x'}\abs{x^\alpha D_x^\beta x'^\gamma D_{x'}^\delta K(x,x')}.
\end{equation*}
For $E\in\Psiiso{-\infty}$ we have $\norm{E}_{\Lin(L^2)}\le\norm{E}_{L^2(\R^{2n})}\le C\norm{E}_{(2n+1)}$.

\begin{lemma}\label{determinants:lem:trilinear}
  If $E_1,E_2\in\Psiiso{-\infty}(\R^n)$ and $T\in\Lin(L^2(\R^n))$, then
  $E_1TE_2\in\Psiiso{-\infty}(\R^n)$ and for each $k$ there are $C$ and $k'$
  such that $\norm{E_1TE_2}_{(k)}\le C\norm{T}_{\Lin(L^2)}\norm{E_1}_{(k')}\norm{E_2}_{(k')}$.
\end{lemma}
\begin{proof}
  The map $x'\mapsto E_2(\cdot,x')$ is smooth into $L^2(\R^n)$, so the kernel
  of $E_1TE_2$ is $K(x,x')=\int E_1(x,w)\,(TE_2(\cdot,x'))(w)\,dw$, and we may
  differentiate under $T$:
  $x^\alpha D_x^\beta x'^\gamma D_{x'}^\delta K(x,x')=\int x^\alpha D_x^\beta E_1(x,w)\,\bigl(T[x'^\gamma D_{x'}^\delta E_2(\cdot,x')]\bigr)(w)\,dw$.
  By Cauchy--Schwarz this is bounded by $\norm{T}$ times
  $\sup_x\norm{x^\alpha D_x^\beta E_1(x,\cdot)}_{L^2}\sup_{x'}\norm{x'^\gamma D_{x'}^\delta E_2(\cdot,x')}_{L^2}$,
  and each factor is bounded by a seminorm $\norm{\cdot}_{(k')}$.
\end{proof}

\begin{lemma}[Finite rank approximation]\label{determinants:lem:approx}
  Let $\pi_N$ be the orthogonal projection onto
  $L_N=\operatorname{span}\{u_\alpha;\abs{\alpha}\le N\}$. For
  $E\in\Psiiso{-\infty}(\R^n)$, $\pi_NE\pi_N\to E$ in $\Psiiso{-\infty}(\R^n)$
  as $N\to\infty$. If $E_y$ is a smooth family, the convergence holds locally
  uniformly in $y$ together with all $y$-derivatives.
\end{lemma}
\begin{proof}
  The functions $u_\alpha\otimes u_\beta$ form an orthonormal basis of
  $L^2(\R^{2n})$ of eigenfunctions of the harmonic oscillator
  $H_{2n}=H_x+H_{x'}$ with eigenvalues $2n+2\abs{\alpha}+2\abs{\beta}$. Thus
  the coefficients $e_{\alpha\beta}=\inner{E}{u_\alpha\otimes u_\beta}$ satisfy
  $(2n+2\abs{\alpha}+2\abs{\beta})^k\abs{e_{\alpha\beta}}=\abs{\inner{H_{2n}^kE}{u_\alpha\otimes u_\beta}}\le\norm{H_{2n}^kE}_{L^2}$,
  a continuous seminorm of $E$. Since the seminorms of $u_\alpha\otimes u_\beta$
  grow only polynomially, $\sum e_{\alpha\beta}u_\alpha\otimes u_\beta$
  converges absolutely to $E$ in $\cS(\R^{2n})$, and its partial sum over
  $\abs{\alpha},\abs{\beta}\le N$ is the kernel of $\pi_NE\pi_N$. For a smooth
  family the same estimates hold for $\pa_y^\gamma E_y$, locally uniformly.
\end{proof}

\begin{lemma}\label{determinants:lem:residual-inverse}
  The set $\Giso{-\infty}(\R^n)$ of $\Id+E$, $E\in\Psiiso{-\infty}(\R^n)$,
  which are invertible on $L^2$ is a group, open in $\Id+\Psiiso{-\infty}(\R^n)$.
  The inverse is $\Id+F$ with
  \begin{equation}\label{determinants:eq:residual-inverse}
    F=-E+E(\Id+E)^{-1}E\in\Psiiso{-\infty}(\R^n),
  \end{equation}
  and $F$ depends continuously on $E$. If $E_y$ is a smooth (holomorphic)
  family with $\Id+E_y$ invertible for all $y$, then $F_y$ is a smooth
  (holomorphic) family.
\end{lemma}
\begin{proof}
  From $(\Id+E)(\Id+F)=\Id=(\Id+F)(\Id+E)$ we get $F=-E-EF$ and
  $F=-E-FE$, so $F=-E+E^2+EFE=-E+E(\Id+E)^{-1}E$, which is residual by
  \cref{determinants:lem:trilinear}. Openness holds because invertibility in
  $\Lin(L^2)$ is open and $\norm{E}_{\Lin(L^2)}\le C\norm{E}_{(2n+1)}$. Inversion
  is smooth (indeed analytic) on the invertible elements of the Banach algebra
  $\Lin(L^2)$, so $y\mapsto(\Id+E_y)^{-1}$ is smooth (holomorphic) in norm; the
  trilinear estimate of \cref{determinants:lem:trilinear} and the product rule
  then show that $F_y$ is smooth (holomorphic), and continuity follows in the
  same way.
\end{proof}

\begin{lemma}[Local inversion]\label{determinants:lem:local-inversion}
  Let $A(z)$ be a holomorphic family of elliptic operators of order $h(z)$ on
  $\Omega$, and suppose $A(p)$ is invertible. Then $A(z)$ is invertible for $z$
  near $p$ and $A(z)^{-1}$ is a holomorphic family of order $-h(z)$ near $p$.
  The same holds for smooth families of fixed order.
\end{lemma}
\begin{proof}
  Let $B(z)$ be a holomorphic parametrix near $p$,
  $B(z)A(z)=\Id+E(z)$, $A(z)B(z)=\Id+E'(z)$. Set
  $B'(z)=B(z)-E(p)A(p)^{-1}$; then $B'(z)A(z)=\Id+E_1(z)$ with
  $E_1(z)=E(z)-E(p)A(p)^{-1}A(z)$ a holomorphic residual family vanishing at
  $p$. By \cref{determinants:lem:residual-inverse}, $\Id+E_1(z)$ is invertible
  near $p$ with holomorphic inverse $\Id+S(z)$, so $L(z)=(\Id+S(z))B'(z)$ is a
  holomorphic family of order $-h(z)$ with $L(z)A(z)=\Id$. In the same way,
  with $B''(z)=B(z)-A(p)^{-1}E'(p)$, we find a right inverse $R(z)$. Then
  $L=LAR=R$ is a two-sided inverse.
\end{proof}

\section{Traces on the residual algebra}\label{determinants:sec:traces}

We start with the simplest invariant. For matrices the trace is
characterized by linearity, the commutator identity and normalization; we
show that the same characterization singles out a unique continuous trace on
$\Psiiso{-\infty}(\R^n)$, given by integrating the kernel over the diagonal.

\subsection*{Matrices}
On the algebra $M(N;\C)$ of $N\times N$ matrices, $\tr A=\sum_iA_{ii}$
satisfies $\tr(AB)=\sum_{i,k}A_{ik}B_{ki}=\tr(BA)$, that is,
\begin{equation*}
  \tr([A,B])=0\quad\forall\,A,B\in M(N;\C),
\end{equation*}
and the normalization $\tr\Id=N$. These properties characterize it.

\begin{lemma}\label{determinants:lem:matrix-trace}
  If $F:M(N;\C)\to\C$ is linear and $F([A,B])=0$ for all $A,B$, then
  $F=F(E_{11})\tr$.
\end{lemma}
\begin{proof}
  Let $E_{jk}$ be the matrix units, $E_{jk}E_{pq}=\delta_{kp}E_{jq}$. For
  $j\ne k$, $E_{jk}=[E_{jj},E_{jk}]$, so $F(E_{jk})=0$; and
  $E_{jj}-E_{11}=[E_{j1},E_{1j}]$, so $F(E_{jj})=F(E_{11})$. Hence
  $F(A)=\sum_{j,k}A_{jk}F(E_{jk})=F(E_{11})\tr A$.
\end{proof}

\subsection*{The trace on residual operators}
For $A\in\Psiiso{-\infty}(\R^n;\C^N)$ with kernel $A(x,x')$ we set
\begin{equation}\label{determinants:eq:trace}
  \Tr(A)=\int_{\R^n}\tr A(x,x)\,dx .
\end{equation}
This is a continuous linear functional on $\Psiiso{-\infty}$. It extends to
operators of sufficiently negative order. If $a\in\Gs{m}(\R^{2n})$ with
$m<-2n$, then $\int\abs{a(x,\xi)}\,d\xi\le C\ang{x}^{m+n}$, so the kernel
$K(x,x')=\int e^{i(x-x')\cdot\xi}a(x,\xi)\,\dbar\xi$ of $\Op(a)$ is continuous,
$\abs{K(x,x)}\le C\ang{x}^{m+n}$ is integrable, and we define
\begin{equation}\label{determinants:eq:trace-symbol}
  \Tr(\Op(a))=\int K(x,x)\,dx=(2\pi)^{-n}\int_{\R^{2n}}\tr a(x,\xi)\,dx\,d\xi,
  \qquad a\in\Gs{m}(\R^{2n}),\ m<-2n.
\end{equation}
This is continuous on $\Psiiso{m}$ for the symbol topology and agrees with
\eqref{determinants:eq:trace} on residual operators. In particular $\Tr$ is
defined on $\Psiisoph{z}$ for $\Re z<-2n$.

\begin{lemma}\label{determinants:lem:trace-property}
  If $A\in\Psiiso{p}(\R^n;\C^N)$, $B\in\Psiiso{q}(\R^n;\C^N)$ and
  $p+q<-2n$, then $\Tr(AB)=\Tr(BA)$. In particular this holds if one of $A$,
  $B$ is residual and the other is arbitrary.
\end{lemma}
\begin{proof}
  For residual $A$ and $B$ this is Fubini's theorem,
  \begin{equation*}
    \Tr(AB)=\iint\tr\bigl(A(x,w)B(w,x)\bigr)\,dw\,dx=\Tr(BA).
  \end{equation*}
  In general, let
  $\chi\in\CIc(\R^{2n})$ equal $1$ near $0$, and let $a_k(Z)=\chi(Z/k)a(Z)$,
  $b_k(Z)=\chi(Z/k)b(Z)$. Then $a-a_k$ is supported in $\abs{Z}\ge k$, where
  $\ang{Z}^{-\eps}\le k^{-\eps}$, and the derivatives of $\chi(Z/k)$ are
  bounded by $C\ang{Z}^{-\abs{\beta}}$ uniformly in $k$; hence $a_k\to a$ in
  $\Gs{p+\eps}$ for every $\eps>0$, and similarly $b_k\to b$ in
  $\Gs{q+\eps}$. Choose $\eps$ with $p+q+2\eps<-2n$. The operators
  $A_k=\Op(a_k)$, $B_k=\Op(b_k)$ are residual, so $\Tr(A_kB_k)=\Tr(B_kA_k)$,
  and by the continuity of composition (see \ref{determinants:I:algebra}) and of
  the trace on $\Psiiso{p+q+2\eps}$ we may let $k\to\infty$.
\end{proof}

\begin{lemma}\label{determinants:lem:finite-trace}
  Let $L\subset\cS(\R^n;\C^N)$ be a finite dimensional subspace and $\Pi$ the
  $L^2$-orthogonal projection onto $L$. Then $\Pi\in\Psiiso{-\infty}$. If
  $T\in\Psiiso{-\infty}$ satisfies $T=\Pi T\Pi$, then $\Tr(T)=\tr(T|_L)$, the
  trace of $T$ as an endomorphism of $L$. More generally
  $\Tr(\Pi T)=\Tr(\Pi T\Pi)=\tr(\Pi T|_L)$ for all $T\in\Psiiso{m}$.
\end{lemma}
\begin{proof}
  If $\varphi_1,\dots,\varphi_l$ is an orthonormal basis of $L$, the kernel of
  $\Pi$ is $\sum_i\varphi_i(x)\overline{\varphi_i(x')}$ (a matrix-valued
  function when $N>1$) and that of $T=\Pi T\Pi$ is
  $\sum_{i,j}t_{ij}\varphi_i(x)\overline{\varphi_j(x')}$, $t_{ij}=\inner{T\varphi_j}{\varphi_i}$.
  Thus $\Tr(T)=\sum_{i,j}t_{ij}\inner{\varphi_i}{\varphi_j}=\sum_it_{ii}$.
  Finally $\Tr(\Pi T(\Id-\Pi))=\Tr(T(\Id-\Pi)\Pi)=0$ by
  \cref{determinants:lem:trace-property}.
\end{proof}

Taking $L=L_N$ in \cref{determinants:lem:finite-trace} and using
\cref{determinants:lem:approx}, we see that
$\Tr(A)=\lim_N\sum_{\abs{\alpha}\le N}\inner{Au_\alpha}{u_\alpha}$, so $\Tr$
coincides with the trace of $A$ as a trace-class operator on $L^2$. The trace
is essentially the only continuous trace on the residual algebra.

\begin{proposition}\label{determinants:prop:trace-unique}
  If $F:\Psiiso{-\infty}(\R^n)\to\C$ is continuous, linear and satisfies
  $F([A,B])=0$ for all $A,B\in\Psiiso{-\infty}(\R^n)$, then $F=c\Tr$ for some
  $c\in\C$. Consequently $F([A,B])=0$ also for all $A\in\Psiiso{-\infty}(\R^n)$
  and $B\in\Psiiso{m}(\R^n)$, $m\in\R$.
\end{proposition}
\begin{proof}
  For a finite dimensional $L\subset\cS(\R^n)$ with orthonormal basis
  $\varphi_1,\dots,\varphi_l$, the operators $P_{ij}$ with kernels
  $\varphi_i(x)\overline{\varphi_j(x')}$ satisfy $P_{ij}P_{kl}=\delta_{jk}P_{il}$,
  so they span a subalgebra of $\Psiiso{-\infty}$ isomorphic to $M(l;\C)$, on
  which $\Tr$ is the matrix trace by \cref{determinants:lem:finite-trace}. By
  \cref{determinants:lem:matrix-trace}, $F=F(P_{11})\Tr$ on this subalgebra.
  Applying this to $L=\operatorname{span}\{\varphi,\psi\}$ for two unit vectors
  $\varphi,\psi\in\cS$ shows that $c=F(\varphi\otimes\overline\varphi)$ is
  independent of $\varphi$. Every finite rank $A\in\Psiiso{-\infty}$ has range,
  and range of $A^*$, in $\cS$, and so lies in the subalgebra associated with
  $L=\Ran A+\Ran A^*$; thus $F=c\Tr$ on finite rank operators, and hence
  everywhere by \cref{determinants:lem:approx} and continuity. The last
  statement then follows from \cref{determinants:lem:trace-property}.
\end{proof}

\begin{discussion}
A useful way to think of $\Psiiso{-\infty}(\R^n)$ is through the Hermite
basis $u_\alpha$. By \cref{determinants:lem:approx}, a residual operator $E$
is the same as an infinite matrix $(e_{\alpha\beta})$ whose entries decay
faster than any power of $|\alpha|+|\beta|$. For such matrices the trace
$\sum_\alpha e_{\alpha\alpha}$ converges absolutely and is the limit of the
traces of the finite truncations $\pi_NE\pi_N$. Most statements in this and
the next section are proved by exactly this reduction to finite matrices,
followed by a limit.
\end{discussion}

\section{The Fredholm determinant}\label{determinants:sec:determinant}

The determinant of $\Id+E$ cannot be defined as a product over a basis,
since there are infinitely many factors. Instead we use the derivative of
$\log\det$, which is a trace, and integrate it along a path.

For matrices, $\det:M(N;\C)\to\C$ is multiplicative, $\det\Id=1$, and
$\GL(N;\C)=\{A;\det A\ne0\}$. Multiplicativity and normalization do not
determine $\det$ (consider $(\det)^k$), but they do together with the
derivative condition $\frac{d}{ds}\det(\Id+sA)|_{s=0}=\tr A$, and for any
smooth curve $A_s$ in $\GL(N;\C)$ one has Jacobi's formula
\begin{equation}\label{determinants:eq:jacobi}
  \frac{d}{ds}\det A_s=\det A_s\,\tr\Bigl(A_s^{-1}\frac{dA_s}{ds}\Bigr)
\end{equation}
(\cref{determinants:ex:matrix-det}). Locally along such a curve this reads
$d\log\det A=\tr(A^{-1}dA)$; the right side is a closed $1$-form on
$\GL(N;\C)$, independent of any choice of logarithm. We use the analogous
$1$-form
\begin{equation*}
  \alpha=\Tr(A^{-1}dA)\ \text{ on }\ \Giso{-\infty}(\R^n),
\end{equation*}
where, at $A=\Id+E$, $dA$ is the identification of tangent vectors with
elements of $\Psiiso{-\infty}$ (so $A^{-1}dA$ takes values in
$\Psiiso{-\infty}$ and $\Tr$ applies). To avoid the differential geometry of
Fr\'echet manifolds we only integrate $\alpha$ along paths. A
\emph{path} in $\Giso{-\infty}$ is a continuous map
$\gamma:[s_0,s_1]\to\Giso{-\infty}$, $\gamma(s)=\Id+E(s)$, which is
piecewise a smooth family in the sense of
\cref{determinants:def:families}, and we set
\begin{equation*}
  \int_\gamma\alpha=\int_{s_0}^{s_1}\Tr\bigl(\gamma(s)^{-1}\gamma'(s)\bigr)\,ds .
\end{equation*}
That $\alpha$ is closed is the content of \cref{determinants:ex:closed}; we
shall not need it, since the following lemma, which compares $\alpha$ with its
finite dimensional analogue, gives path independence directly.

\begin{lemma}\label{determinants:lem:connected}
  $\Giso{-\infty}(\R^n)$ is path connected: if $\Id+E\in\Giso{-\infty}$ there is
  a path in $\Giso{-\infty}$ from $\Id$ to $\Id+E$.
\end{lemma}
\begin{proof}
  $E$ is compact on $L^2$, so its non-zero spectrum consists of isolated
  eigenvalues accumulating only at $0$. For $c\ne0$, $\Id+cE$ is invertible
  unless $-1/c$ is an eigenvalue, so the set $\Lambda$ of $c$ with
  $\abs{c}\le2$ and $\Id+cE$ not invertible is finite, and $0,1\notin\Lambda$.
  If $c:[0,1]\to\{\abs{c}\le2\}\setminus\Lambda$ is a smooth path from $0$ to
  $1$, then $\Id+c(s)E$ is the required path.
\end{proof}

\begin{figure}[htbp]
\centering
\begin{tikzpicture}[>=Latex]
  % (a) the c-plane
  \begin{scope}
    \fill[gray!10] (0,0) circle (2);
    \draw (0,0) circle (2);
    \draw[->,gray] (-2.3,0) -- (2.4,0) node[right,black] {\small $\operatorname{Re}c$};
    \draw[->,gray] (0,-2.3) -- (0,2.4) node[above,black] {\small $\operatorname{Im}c$};
    \foreach \p in {(0.5,0),(-0.9,0.8),(1.3,-0.9),(-1.2,-1.1),(0.7,1.2)}
      \node at \p {\small $\times$};
    \fill (0,0) circle (1.6pt) node[below left] {\small $0$};
    \fill (1,0) circle (1.6pt) node[below right] {\small $1$};
    \draw[thick,->] (0,0) .. controls (0.15,0.6) and (0.8,0.55) .. (1,0);
    \node at (0.12,0.72) {\small $c(s)$};
    \node[above right] at (1.35,1.45) {\small $|c|\le2$};
    \node at (0,-2.75) {\small (a) $\Lambda$ ($\times$) is finite};
  \end{scope}
  % (b) the residual group
  \begin{scope}[xshift=6.6cm]
    \fill[gray!10] (-2.4,-2) rectangle (2.4,2);
    \draw (-2.4,-2) rectangle (2.4,2);
    \draw[very thick,gray] (-2.4,0.6) .. controls (-1,1.4) and (0,-0.2) .. (0.8,2);
    \draw[very thick,gray] (0.3,-2) .. controls (0.6,-0.6) and (1.6,0.2) .. (2.4,-0.4);
    \node[gray!60!black] at (1.45,-1.45) {\small $\detFr=0$};
    \fill (-1.5,-1.2) circle (1.6pt) node[below] {\small $\Id$};
    \fill (1.6,1.2) circle (1.6pt) node[right] {\small $\Id+E$};
    \draw[thick,->] (-1.5,-1.2) .. controls (-0.5,-0.6) and (0.2,0.2) .. (1.6,1.2);
    \node at (0.15,-0.55) {\small $\gamma$};
    \node at (0,-2.75) {\small (b) a path in $\Giso{-\infty}$};
  \end{scope}
\end{tikzpicture}
\caption[Paths in the residual group]{(a) The proof of
\cref{determinants:lem:connected}: the values of $c$ for which $\Id+cE$ is
not invertible form a finite set in $\{|c|\le2\}$, which a path from $0$ to
$1$ can avoid. (b) Schematic picture of $\Id+\Psiiso{-\infty}(\R^n)$: the
non-invertible operators form a (singular) hypersurface, the zero set of
$\detFr$, and $\detFr(\Id+E)$ is defined by integrating
$\Tr(\gamma^{-1}\gamma')$ along any path $\gamma$ from $\Id$ that avoids it.}
\label{determinants:fig:det-path}
\end{figure}

\begin{lemma}\label{determinants:lem:path-limit}
  Let $\gamma(s)=\Id+E(s)$, $s\in[0,1]$, be a path in $\Giso{-\infty}(\R^n)$
  with $E(0)=0$. Then
  \begin{equation*}
    \exp\Bigl(\int_\gamma\alpha\Bigr)=\lim_{N\to\infty}
    \det\bigl(\Id+\pi_NE(1)|_{L_N}\bigr),
  \end{equation*}
  where $\pi_NE(1)|_{L_N}$ is regarded as an endomorphism of the finite
  dimensional space $L_N$. In particular the left side depends only on
  $\gamma(1)$.
\end{lemma}
\begin{proof}
  Set $E_N(s)=\pi_NE(s)\pi_N$ and $\gamma_N=\Id+E_N$. By
  \cref{determinants:lem:approx}, $E_N\to E$ and $E_N'\to E'$ in
  $\Psiiso{-\infty}$, uniformly on each interval of smoothness. Let
  $M=\sup_s\norm{\gamma(s)^{-1}}_{\Lin(L^2)}<\infty$. For $N$ large,
  $\norm{E_N(s)-E(s)}_{\Lin(L^2)}<1/(2M)$ for all $s$, so $\gamma_N(s)$ is
  invertible with $\gamma_N(s)^{-1}\to\gamma(s)^{-1}$ uniformly in norm. By
  \eqref{determinants:eq:residual-inverse} and
  \cref{determinants:lem:trilinear}, $\gamma_N^{-1}-\Id\to\gamma^{-1}-\Id$
  uniformly in $\Psiiso{-\infty}$, hence
  $\gamma_N^{-1}\gamma_N'\to\gamma^{-1}\gamma'$ uniformly in
  $\Psiiso{-\infty}$, and, $\Tr$ being continuous,
  $\int_{\gamma_N}\alpha\to\int_\gamma\alpha$.

  Now $\gamma_N(s)$ preserves the decomposition $L^2=L_N\oplus L_N^\perp$,
  acting as $g_N(s)=\Id+\pi_NE(s)|_{L_N}$ on $L_N$ and as the identity on
  $L_N^\perp$. Hence $\gamma_N^{-1}\gamma_N'=\pi_N\gamma_N^{-1}\gamma_N'\pi_N$
  acts as $g_N^{-1}g_N'$ on $L_N$, and by
  \cref{determinants:lem:finite-trace} and Jacobi's formula
  \eqref{determinants:eq:jacobi},
  \begin{equation*}
    \int_{\gamma_N}\alpha=\int_0^1\tr\bigl(g_N^{-1}g_N'\bigr)\,ds,\qquad
    \exp\int_{\gamma_N}\alpha=\frac{\det g_N(1)}{\det g_N(0)}=\det g_N(1).
  \end{equation*}
  Letting $N\to\infty$ proves the lemma.
\end{proof}

\begin{definition}\label{determinants:def:det}
  The \emph{Fredholm determinant} $\detFr:\Id+\Psiiso{-\infty}(\R^n;\C^N)\to\C$ is
  \begin{equation*}
    \detFr(\Id+E)=\begin{cases}
      \exp\bigl(\int_\gamma\alpha\bigr) & \text{if }\Id+E\in\Giso{-\infty},\\
      0&\text{otherwise,}
    \end{cases}
  \end{equation*}
  where $\gamma$ is any path in $\Giso{-\infty}$ from $\Id$ to $\Id+E$.
\end{definition}

By \cref{determinants:lem:connected,determinants:lem:path-limit} this is well
defined (see \cref{determinants:fig:det-path}).

\begin{question}
Why not simply define $\detFr(\Id+E)=\prod_j(1+\lambda_j)$, the product
over the eigenvalues $\lambda_j$ of $E$?
\end{question}
\begin{answer}
This is in fact correct for trace class $E$ (a theorem of Lidskii), but
proving it requires a good deal of spectral theory of non-self-adjoint
compact operators, and the product formula is awkward for proving
multiplicativity or holomorphy. The path definition uses only the trace
property and finite-dimensional determinants, and the limit formula of
\cref{determinants:lem:path-limit} shows directly that it agrees with the
determinants of the finite truncations $\pi_NE\pi_N$.
\end{answer}

\noindent Concatenating a path from $\Id$ to $G_1$ with a path $\delta$ in
$\Giso{-\infty}$ from $G_1$ to $G_2$ gives
\begin{equation}\label{determinants:eq:det-variation}
  \detFr(G_2)=\detFr(G_1)\exp\Bigl(\int_\delta\alpha\Bigr).
\end{equation}
In particular, if $G_y$ is a smooth family in $\Giso{-\infty}$, then
$\detFr(G_y)$ is smooth and $d\log\detFr(G_y)=\Tr(G_y^{-1}dG_y)$: applying
\eqref{determinants:eq:det-variation} along curves shows that the directional
derivatives are $\detFr(G)\Tr(G^{-1}\pa G)$, which are continuous, and
iterating gives smoothness. For a holomorphic family the Cauchy--Riemann
equations follow in the same way, so $\detFr(G(z))$ is holomorphic. Also
$\detFr$ is continuous on $\Giso{-\infty}$, since for $G'$ close to $G$ the
segment from $G$ to $G'$ lies in $\Giso{-\infty}$.

\begin{proposition}\label{determinants:prop:det-basic}
  The Fredholm determinant has the following properties.
  \begin{enumerate}[label=(\roman*)]
  \item\label{determinants:it:mult} $\detFr(AB)=\detFr(A)\detFr(B)$ for
    $A,B\in\Id+\Psiiso{-\infty}$, and $\detFr(\Id)=1$.
  \item $\detFr(A)\ne0$ if and only if $A$ is invertible.
  \item\label{determinants:it:explog} $\detFr(\exp E)=\exp\Tr(E)$ for all
    $E\in\Psiiso{-\infty}$; if $\norm{E}_{\Lin(L^2)}<1$, then
    $\log(\Id+E)=\sum_{k\ge1}(-1)^{k+1}E^k/k$ converges in $\Psiiso{-\infty}$
    and $\detFr(\Id+E)=\exp\Tr\log(\Id+E)$.
  \item\label{determinants:it:conj} If $A\in\Psiisoph{m}(\R^n;\C^N)$ is
    invertible (with inverse in $\Psiisoph{-m}$) and $B\in\Id+\Psiiso{-\infty}$,
    then $ABA^{-1}\in\Id+\Psiiso{-\infty}$ and $\detFr(ABA^{-1})=\detFr(B)$.
  \end{enumerate}
\end{proposition}
\begin{proof}
  (i) If $A$ and $B$ are invertible, let $\gamma_1$, $\gamma_2$ be paths from
  $\Id$ to $A$ and $B$; then $\gamma=\gamma_1\gamma_2$ is a path from $\Id$ to
  $AB$ and
  $\gamma^{-1}\gamma'=\gamma_2^{-1}(\gamma_1^{-1}\gamma_1')\gamma_2+\gamma_2^{-1}\gamma_2'$.
  By \cref{determinants:lem:trace-property},
  $\Tr(\gamma_2^{-1}X\gamma_2)=\Tr(X\gamma_2\gamma_2^{-1})=\Tr(X)$ for residual
  $X$, so $\int_\gamma\alpha=\int_{\gamma_1}\alpha+\int_{\gamma_2}\alpha$. If
  $AB$ is invertible then $A$ is surjective and $B$ is injective, so both are
  invertible since they have index zero \ref{determinants:I:residual}. Hence
  if one of $A,B$ is not invertible, neither is $AB$, and both sides vanish.

  (ii) is immediate from the definition.

  (iii) The path $s\mapsto\exp(sE)$ has $\gamma^{-1}\gamma'=E$. If
  $\norm{E}<1$, the path $\Id+sE$, $s\in[0,1]$, lies in $\Giso{-\infty}$, and
  $(\Id+sE)^{-1}E=\sum_{k\ge0}(-s)^kE^{k+1}$, the series converging in
  $\Psiiso{-\infty}$ uniformly for $s\in[0,1]$ because, by
  \cref{determinants:lem:trilinear},
  $\norm{E^{k+2}}_{(j)}=\norm{E\,E^{k}E}_{(j)}\le C\norm{E}_{\Lin(L^2)}^k\norm{E}_{(j')}^2$.
  Integrating term by term gives
  $\int_\gamma\alpha=\sum_{k\ge0}(-1)^k\Tr(E^{k+1})/(k+1)=\Tr\log(\Id+E)$.

  (iv) $ABA^{-1}-\Id=A(B-\Id)A^{-1}$ is residual and $ABA^{-1}$ is invertible
  if and only if $B$ is. If $\gamma$ is a path from $\Id$ to $B$, then
  $A\gamma A^{-1}$ is a path from $\Id$ to $ABA^{-1}$ and
  $\Tr\bigl(A\gamma^{-1}\gamma'A^{-1}\bigr)=\Tr(\gamma^{-1}\gamma')$ by
  \cref{determinants:lem:trace-property}.
\end{proof}

\begin{lemma}[Finite dimensional reduction]\label{determinants:lem:reduction}
  Let $\Pi$ be the orthogonal projection onto a finite dimensional subspace
  $L\subset\cS(\R^n;\C^N)$, and let $R\in\Psiiso{-\infty}$ satisfy $R=R\Pi$.
  Then $\Id+R$ is invertible if and only if $\Id_L+\Pi R|_L$ is invertible, and
  \begin{equation*}
    \detFr(\Id+R)=\det\bigl(\Id_L+\Pi R|_L\bigr).
  \end{equation*}
  The same holds if $R=\Pi R$.
\end{lemma}
\begin{proof}
  With respect to $L^2=L\oplus L^\perp$ (columns corresponding to the
  domain), $\Id+cR=\begin{pmatrix}\Id+cr&0\\c\rho&\Id\end{pmatrix}$ with
  $r=\Pi R|_L$ and $\rho=(\Id-\Pi)R|_L$. This is invertible if and only if
  $q(c)=\det(\Id_L+cr)\ne0$, which proves the first statement for $c=1$. If
  $q(1)=0$ both sides vanish. Otherwise, since the polynomial $q$ has
  $q(0)=1$, it has finitely many zeros, and we may choose a smooth path $c(s)$
  from $0$ to $1$ avoiding them; $\gamma(s)=\Id+c(s)R$ is then a path in
  $\Giso{-\infty}$. From the triangular form,
  $T=\gamma^{-1}\gamma'=\begin{pmatrix}(\Id+cr)^{-1}c'r&0\\ *&0\end{pmatrix}$,
  so $T=T\Pi$ and, by \cref{determinants:lem:finite-trace} and
  \cref{determinants:lem:trace-property},
  $\Tr(T)=\Tr(\Pi T\Pi)+\Tr((\Id-\Pi)T\Pi)=\tr\bigl((\Id+cr)^{-1}c'r\bigr)+\Tr(\Pi(\Id-\Pi)T)$.
  The last term vanishes and by Jacobi's formula the first is
  $\frac{d}{ds}\log q(c(s))$, so $\exp\int_\gamma\alpha=q(1)/q(0)=\det(\Id_L+r)$.
  The case $R=\Pi R$ follows by taking adjoints, or by the same argument with
  upper triangular matrices.
\end{proof}

In particular, for finite rank residual operators the Fredholm determinant is
an ordinary determinant, and $\detFr(\Id+\pi_NE\pi_N)=\det(\Id+\pi_NE|_{L_N})$.

\begin{theorem}\label{determinants:thm:det-entire}
  The Fredholm determinant is continuous on $\Id+\Psiiso{-\infty}(\R^n;\C^N)$.
  If $E_y$ is a smooth (resp.\ holomorphic) family in $\Psiiso{-\infty}$, then
  $\detFr(\Id+E_y)$ is a smooth (resp.\ holomorphic) function. Moreover, for
  every $E\in\Psiiso{-\infty}$,
  \begin{equation}\label{determinants:eq:det-limit}
    \detFr(\Id+E)=\lim_{N\to\infty}\det\bigl(\Id+\pi_NE|_{L_N}\bigr).
  \end{equation}
\end{theorem}
\begin{proof}
  Near invertible points this was shown above. Let $\Id+E_0$ be
  non-invertible, let $\Pi_0$ be the projection onto $\Nul(\Id+E_0)\subset\cS$,
  and let $F$ be as in \ref{determinants:I:residual}, so that
  $G_0=\Id+E_0+F$ is invertible and $F=F\Pi_0$. For $E$ near $E_0$ the operator
  $G_0+E-E_0$ is invertible; let $K(E)=(G_0+E-E_0)^{-1}$, which lies in
  $\Giso{-\infty}$ and depends continuously on $E$
  (\cref{determinants:lem:residual-inverse}). Then
  \begin{equation*}
    \Id+E=K(E)^{-1}\bigl(\Id-K(E)F\bigr),\qquad
    \detFr(\Id+E)=\detFr(K(E))^{-1}\det\bigl(\Id-\Pi_0K(E)F|_{L}\bigr),
  \end{equation*}
  using \cref{determinants:prop:det-basic}\ref{determinants:it:mult} and
  \cref{determinants:lem:reduction} with $L=\Ran\Pi_0$ and $R=-K(E)F=R\Pi_0$.
  The first factor is continuous and non-vanishing, and the second is a
  polynomial in the matrix entries $\inner{K(E)Fe_j}{e_i}$ for an orthonormal
  basis $e_i$ of $L$, which are continuous in $E$. If $E=E_y$ is a smooth or
  holomorphic family, both factors are smooth or holomorphic. Finally
  \eqref{determinants:eq:det-limit} follows from
  $\detFr(\Id+\pi_NE\pi_N)=\det(\Id+\pi_NE|_{L_N})$, the convergence
  $\pi_NE\pi_N\to E$ and continuity.
\end{proof}

\begin{remark}\label{determinants:rem:det-entire}
  Thus $\detFr$ is an entire function on $\Id+\Psiiso{-\infty}$: holomorphic
  along every holomorphic family, and in particular $c\mapsto\detFr(\Id+cE)$
  is entire. In view of \cref{determinants:prop:det-basic}, $\Giso{-\infty}$
  is the complement of the (singular) hypersurface
  $\{\Id+E;\ \detFr(\Id+E)=0\}$. The trace and the determinant extend further,
  to operators of order less than $-2n$; see \cref{determinants:ex:extend}.
\end{remark}

\section{Index and invertible perturbations}\label{determinants:sec:index}

For an elliptic $A\in\Psiisoph{m}(\R^n;\C^N)$ the null spaces of $A$ and $A^*$
on $\cSp$ are finite dimensional subspaces of $\cS$, and the \emph{index} is
\begin{equation*}
  \ind(A)=\dim\Nul(A)-\dim\Nul(A^*).
\end{equation*}
For $m=0$ this is the index of $A$ as a Fredholm operator on $L^2$, since
$\Ran(A)$ is closed with orthocomplement $\Nul(A^*)$ there.

\begin{proposition}\label{determinants:prop:index-trace}
  If $A\in\Psiisoph{m}(\R^n;\C^N)$ is elliptic and $B$ is any parametrix, then
  $[A,B]=(AB-\Id)-(BA-\Id)\in\Psiiso{-\infty}$ and
  \begin{equation}\label{determinants:eq:index-trace}
    \ind(A)=\Tr([A,B]).
  \end{equation}
  Consequently $\ind(A+E)=\ind(A)$ for $E\in\Psiiso{-\infty}$, and if $A(z)$
  is a holomorphic (or smooth) family of elliptic operators on a connected set,
  then $\ind A(z)$ is constant.
\end{proposition}
\begin{proof}
  For the generalized inverse, \eqref{determinants:eq:gen-inverse} and
  \cref{determinants:lem:finite-trace} give
  $\ind(A)=\Tr\Pi_0-\Tr\Pi_1=\Tr(\Id-BA)-\Tr(\Id-AB)=\Tr([A,B])$. Any other
  parametrix is $B'=B+E'$ with $E'$ residual, and
  $\Tr([A,B'])-\Tr([A,B])=\Tr([A,E'])=0$ by
  \cref{determinants:lem:trace-property}. A parametrix of $A$ is also one of
  $A+E$, and $\Tr([E,B])=0$. For a family, choose locally a holomorphic
  parametrix $B(z)$ \ref{determinants:I:families}; then $[A(z),B(z)]$ is a
  holomorphic residual family, so $\ind A(z)$ is continuous, hence locally
  constant, hence constant.
\end{proof}

\begin{discussion}
For matrices $\Tr([A,B])=0$, so formula \eqref{determinants:eq:index-trace}
recovers the fact that a square matrix has index zero. In infinite
dimensions neither $A$ nor $B$ need be of trace class, and $[A,B]$ can have
non-zero trace even though it is a commutator. The index measures exactly
this failure of the trace property. The rest of the chapter makes this
precise: once the trace is extended (as the regularized trace) to all of
$\Psiisoph{\Z}$, its failure to vanish on commutators is a local
expression in the symbols, and so is the index
(\cref{determinants:thm:index-formula}).
\end{discussion}

\begin{proposition}\label{determinants:prop:invertible-perturbation}
  Let $A\in\Psiisoph{m}(\R^n;\C^N)$ be elliptic with $\ind(A)=0$. Then there is
  a finite rank $F\in\Psiiso{-\infty}(\R^n;\C^N)$ such that $A+sF$ is
  invertible for every $s\ne0$; its inverse lies in $\Psiisoph{-m}$.
\end{proposition}
\begin{proof}
  Let $e_1,\dots,e_p$ and $f_1,\dots,f_p$ be orthonormal bases of $\Nul(A)$ and
  $\Nul(A^*)$, both in $\cS$. Let $F$ and $G$ be the operators with kernels
  $\sum_jf_j(x)\overline{e_j(x')}$ and $\sum_je_j(x)\overline{f_j(x')}$, so
  $F=\Pi_1F=F\Pi_0$, $G=\Pi_0G=G\Pi_1$, $FG=\Pi_1$, $GF=\Pi_0$. With $B$ the
  generalized inverse, normalized as in \ref{determinants:I:elliptic}, and
  using $A\Pi_0=0$, $\Pi_1A=0$, $\Pi_0B=0$, $B\Pi_1=0$,
  \begin{gather*}
    (A+sF)(B+s^{-1}G)=AB+s^{-1}A\Pi_0G+sF\Pi_0B+FG=\Id-\Pi_1+\Pi_1=\Id,\\
    (B+s^{-1}G)(A+sF)=BA+sB\Pi_1F+s^{-1}G\Pi_1A+GF=\Id .
  \end{gather*}
  Thus $A+sF$ is invertible with inverse $B+s^{-1}G\in\Psiisoph{-m}$.
\end{proof}

Write $\Ell^0_{\mathrm{iso}}(\R^n;\C^N)\subset\Psiisoph{0}(\R^n;\C^N)$ for the
elliptic elements, $\Ell^0_{\mathrm{iso},k}$ for those of index $k$, and
$\Giso{0}(\R^n;\C^N)$ for the group of invertible elements (on $L^2$, or
equivalently in $\Psiisoph{0}$). Then
$\Giso{0}\subset\Ell^0_{\mathrm{iso},0}\subset\Ell^0_{\mathrm{iso}}\subset\Psiisoph{0}$.

\begin{corollary}\label{determinants:cor:dense}
  $\Giso{0}(\R^n;\C^N)$ is dense in $\Ell^0_{\mathrm{iso},0}(\R^n;\C^N)$ for
  the topology of $\Psiisoph{0}$: indeed $A+sF\to A$ as $s\to0$.
\end{corollary}

By \cref{determinants:prop:index-trace} the index is defined on the quotient.
Let
\begin{equation*}
  \mathcal A^m=\Psiisoph{m}(\R^n;\C^N)/\Psiiso{-\infty}(\R^n;\C^N);
\end{equation*}
then $\mathcal A^{\Z}=\bigcup_m\mathcal A^m$ is an algebra, $A$ is elliptic if and
only if its class $a=[A]$ is invertible in $\mathcal A^{\Z}$ (the inverse
being the class of a parametrix), and
\begin{equation*}
  \ind:\{a\in\mathcal A^0\text{ invertible}\}\to\Z,\qquad \ind(a)=\ind(A)\ \text{ if }a=[A],
\end{equation*}
is well defined. We return to it in \cref{determinants:sec:index-formula}.

\section{Analytic Fredholm theory}\label{determinants:sec:fredholm}

The determinant turns questions about invertibility of operators into
questions about zeros of a scalar function. For a holomorphic family
$\Id+A(z)$ that function is holomorphic, so its zeros are isolated, and the
inverse can only blow up at isolated points.

A family $M(z)$ with values in $\Psiiso{-\infty}$ is \emph{meromorphic} near
$p$ if $M(z)=\sum_{k=1}^K(z-p)^{-k}M_{-k}+M_0(z)$ with $M_{-k}\in\Psiiso{-\infty}$
and $M_0$ a holomorphic family; we say it has \emph{finite rank singular part}
if the $M_{-k}$ have finite rank.

\begin{lemma}[Fredholm alternative]\label{determinants:lem:fredholm-alt}
  Let $\Omega\subset\C$ be open and connected and let $A(z)$ be a holomorphic
  family in $\Psiiso{-\infty}(\R^n;\C^N)$. Then either $\Id+A(z)$ is invertible
  for no $z\in\Omega$, or it is invertible outside a discrete set
  $D\subset\Omega$; in the latter case $(\Id+A(z))^{-1}=\Id+B(z)$ where $B$ is
  meromorphic in $\Psiiso{-\infty}$ with finite rank singular parts.
\end{lemma}
\begin{proof}
  By \cref{determinants:thm:det-entire}, $f(z)=\detFr(\Id+A(z))$ is
  holomorphic on $\Omega$, so either $f\equiv0$ or its zero set $D$ is
  discrete; and $\Id+A(z)$ is invertible exactly when $f(z)\ne0$. Away from
  $D$, $B(z)$ is holomorphic by \cref{determinants:lem:residual-inverse}.
  Near $p\in D$ we proceed as in the proof of
  \cref{determinants:thm:det-entire}: with $L=\Nul(\Id+A(p))$, $\Pi_0$ the
  projection onto $L$ and $G_0=\Id+A(p)+F$ invertible, $F=F\Pi_0$ of finite
  rank, the family $K(z)=(G_0+A(z)-A(p))^{-1}$ is holomorphic in
  $\Giso{-\infty}$ near $p$ and $\Id+A(z)=K(z)^{-1}(\Id+R(z))$ with
  $R(z)=-K(z)F=R(z)\Pi_0$. The matrix $\Id_L+\Pi_0R(z)|_L$ has determinant
  $\detFr(\Id+R(z))=f(z)\detFr(K(z))\not\equiv0$, so its inverse $M(z)$ is a
  meromorphic matrix (Cramer's rule). A direct computation, using
  $R=R\Pi_0$ and $(\Id_L+\Pi_0R\Pi_0)M=\Id_L$, shows
  $(\Id+R)(\Id-RM\Pi_0)=\Id+R-R\Pi_0=\Id$, so
  \begin{equation*}
    (\Id+A(z))^{-1}=K(z)-\sum_{i,j}m_{ij}(z)\,R(z)P_{ij}K(z),
  \end{equation*}
  where $P_{ij}$ are the matrix units with respect to an orthonormal basis of
  $L$ and $m_{ij}$ the entries of $M$. Each $R(z)P_{ij}K(z)$ is a holomorphic
  family of residual operators of rank at most one, whose Taylor coefficients
  therefore have finite rank; multiplying by the meromorphic scalars $m_{ij}$
  gives a singular part of finite rank.
\end{proof}

This result for smoothing operators, which really follows from the
corresponding result for matrices, implies the analogous statement for
holomorphic families of elliptic operators. The idea of the proof is to
perturb by a fixed finite rank operator $F$ so as to make $A(z)+F$
invertible near a point; then $A(z)=(A(z)+F)(\Id-C(z)F)$, and the second
factor is a residual perturbation of the identity, to which
\cref{determinants:lem:fredholm-alt} applies.

\begin{proposition}\label{determinants:prop:analytic-fredholm}
  Let $\Omega\subset\C$ be open and connected, $h$ holomorphic on $\Omega$, and
  $A(z)$ a holomorphic family of elliptic operators of order $h(z)$. If $A(z_0)$
  is invertible for some $z_0\in\Omega$, then $A(z)$ is invertible outside a
  discrete set $D\subset\Omega$, and near each $p\in\Omega$
  \begin{equation*}
    A(z)^{-1}=B(z)+H(z),
  \end{equation*}
  with $B(z)$ a holomorphic family of order $-h(z)$ and $H(z)$ meromorphic in
  $\Psiiso{-\infty}$ with finite rank singular part ($H=0$ if $p\notin D$).
\end{proposition}
\begin{proof}
  By \cref{determinants:prop:index-trace}, $\ind A(z)=\ind A(z_0)=0$ for all
  $z$. Fix $p\in\Omega$. By \cref{determinants:prop:invertible-perturbation}
  there is a finite rank residual $F$ with $A(p)+F$ invertible, so by
  \cref{determinants:lem:local-inversion} $A(z)+F$ is invertible on a disc
  $U_p$ around $p$, with inverse $C(z)$ a holomorphic family of order $-h(z)$.
  Then $A(z)=(A(z)+F)(\Id-C(z)F)$, so on $U_p$, $A(z)$ is invertible exactly
  where $g_p(z)=\detFr(\Id-C(z)F)\ne0$.

  Let $W$ be the set of $p\in\Omega$ such that $A(z)$ is invertible at points
  arbitrarily close to $p$. If $p\in W$ then $g_p\not\equiv0$ on the disc
  $U_p$, so its zeros are isolated, and every point of $U_p$ is in $W$; thus
  $W$ is open. Its complement is the interior of the set of non-invertible
  points, also open. Since $z_0\in W$ and $\Omega$ is connected, $W=\Omega$,
  and the non-invertible points form a discrete set.

  Near $p$, \cref{determinants:lem:fredholm-alt} gives
  $(\Id-C(z)F)^{-1}=\Id+M(z)$ with $M=\sum_{k=1}^K(z-p)^{-k}M_{-k}+M_0(z)$ as
  above, so $A(z)^{-1}=(\Id+M(z))C(z)$. Set $B=C+M_0C$, a holomorphic family of
  order $-h(z)$, and $H=\sum_k(z-p)^{-k}M_{-k}C(z)$. Each $M_{-k}C(z)$ is a
  holomorphic residual family with range in the finite dimensional space
  $\Ran M_{-k}$, so the singular part of $H$ has finite rank.
\end{proof}

\section{Resolvent and spectrum}\label{determinants:sec:resolvent}

One direct application of analytic Fredholm theory is to the resolvent of an
elliptic operator of positive order. For simplicity we take the order to be a
positive integer $m$, so that $A-\lambda\in\Psiisoph{m}$ for $\lambda\in\C$; the
case of non-integral order needs only notational changes, replacing
$\Psiisoph{m}$ by the class of operators with symbols in $\Gs{m}$ having an
expansion in terms of homogeneous functions of degrees $m-j$ and $-j$.

\begin{proposition}\label{determinants:prop:resolvent}
  Let $A\in\Psiisoph{m}(\R^n;\C^N)$, $m\in\N$, be elliptic, and suppose that for
  some $\lambda'\in\C$ both $A-\lambda'$ and $A^*-\overline{\lambda'}$ have
  trivial null space. Then the set $\Spec(A)$ of $\lambda\in\C$ for which
  $A-\lambda$ is not invertible is discrete, and $(A-\lambda)^{-1}$ is a
  meromorphic family in $\Psiisoph{-m}$ with finite rank residual singular
  parts. At $\tilde\lambda\in\Spec(A)$ the generalized eigenspace
  \begin{equation*}
    E_{\tilde\lambda}=\{u\in\cSp(\R^n;\C^N);\ (A-\tilde\lambda)^pu=0\text{ for some }p\in\N\}
  \end{equation*}
  is a finite dimensional subspace of $\cS$, and
  \begin{equation*}
    (A-\lambda)^{-1}=-\sum_{j=1}^k(\lambda-\tilde\lambda)^{-j}(A-\tilde\lambda)^{j-1}P+(\text{holomorphic near }\tilde\lambda),
  \end{equation*}
  where $P$ is a projection (not necessarily orthogonal) onto $E_{\tilde\lambda}$
  commuting with $A$, and $k$ is the least integer with $(A-\tilde\lambda)^kP=0$.
\end{proposition}
\begin{proof}
  Since $m>0$, $A-\lambda$ is, for every $\lambda$, elliptic with the same
  principal symbol as $A$, and $\lambda\mapsto A-\lambda$ is an entire family
  of order $m$. The hypothesis says $\ind(A-\lambda')=0$ and $A-\lambda'$ is
  invertible, so \cref{determinants:prop:analytic-fredholm} applies with
  $\Omega=\C$. Replacing $A$ by $A-\tilde\lambda$ we may take $\tilde\lambda=0$
  and write $(A-\lambda)^{-1}=\sum_{j\ge-k}\lambda^jS_j$, with $S_{-1},\dots,S_{-k}$
  finite rank residual. Comparing coefficients in
  $(A-\lambda)(A-\lambda)^{-1}=\Id=(A-\lambda)^{-1}(A-\lambda)$ gives
  $AS_j-S_{j-1}=\delta_{j0}\Id=S_jA-S_{j-1}$, so with $R_j=S_{-j}$,
  \begin{equation*}
    R_{j+1}=AR_j=R_jA\ (1\le j\le k),\quad R_{k+1}=0,\qquad\text{hence } R_j=A^{j-1}R_1,\ A^kR_1=0.
  \end{equation*}
  Thus $\Ran R_1\subset E_0$ and $R_1$ commutes with $A$. Conversely, if $A^pu=0$,
  then $(A-\lambda)\bigl(-\sum_{j=0}^{p-1}\lambda^{-j-1}A^ju\bigr)=u$, so
  $(A-\lambda)^{-1}u=-\sum_{j<p}\lambda^{-j-1}A^ju$ and comparing the
  coefficients of $\lambda^{-1}$, $R_1u=-u$. Hence $P=-R_1$ is a projection
  with range exactly $E_0$, which is therefore finite dimensional and contained
  in $\cS$.
\end{proof}

Note also that $\Ran R_k\subset\Ran R_{k-1}\subset\dots\subset\Ran R_1$ and
$\Nul R_k\supset\dots\supset\Nul R_1$; precisely, $u\in\Ran R_p$, $p\ge1$, if
and only if $u=A^{p-1}u_1$ for some $u_1\in\Ran R_1$.

\begin{discussion}
In the notation of the proof, $P=-R_1$ is minus the residue of the
resolvent at $\tilde\lambda$, so by the residue theorem
\begin{equation*}
  P=-\frac1{2\pi i}\oint_{|\lambda-\tilde\lambda|=\eps}(A-\lambda)^{-1}\,d\lambda
\end{equation*}
for small $\eps>0$ (\cref{determinants:fig:resolvent}). This is the usual
Riesz projection of spectral theory. Here no self-adjointness is assumed, so
eigenvalues may be complex and the eigenspaces need not be orthogonal; the
analytic Fredholm theorem still guarantees that the spectrum is discrete
and that each generalized eigenspace is finite dimensional.
\end{discussion}

\begin{figure}[htbp]
\centering
\begin{tikzpicture}[>=Latex]
  \draw[->,gray] (-1,0) -- (8,0) node[right,black] {\small $\operatorname{Re}\lambda$};
  \draw[->,gray] (0,-1.6) -- (0,2.0) node[above,black] {\small $\operatorname{Im}\lambda$};
  \foreach \p in {(1.2,0.3),(2.0,-0.5),(2.9,0.9),(3.6,-0.2),(4.6,0.6),(5.9,-0.9),(6.4,0.4),(6.9,-0.3),(7.4,1.0)}
    \fill \p circle (1.8pt);
  \fill (3.6,-0.2) circle (1.8pt);
  \node[below] at (3.6,-0.25) {\small $\tilde\lambda$};
  \draw[thick,->] (3.6,-0.2) ++(0:0.55) arc (0:350:0.55);
  \node at (4.3,-1.2) {\small $|\lambda-\tilde\lambda|=\eps$};
  \node at (-0.45,1.35) {\small $\times$};
  \node[above] at (-0.45,1.4) {\small $\lambda'$};
  \node[font=\small] at (5.6,1.65) {$\Spec(A)$ discrete};
\end{tikzpicture}
\caption[Spectrum and resolvent]{The spectrum of an elliptic $A$ of
positive order satisfying the hypothesis of
\cref{determinants:prop:resolvent} is a discrete set (dots) in $\C$; at
$\lambda'$ ($\times$) both $A-\lambda'$ and its adjoint are injective. The
resolvent $(A-\lambda)^{-1}$ is meromorphic, and integrating it around a
small circle about $\tilde\lambda\in\Spec(A)$ gives $-2\pi i$ times the
projection $P$ onto the generalized eigenspace $E_{\tilde\lambda}$.}
\label{determinants:fig:resolvent}
\end{figure}

\section{The residue trace}\label{determinants:sec:residue}

We have shown that there is a unique trace on the residual algebra. Following
ideas of Seeley \cite{Seeley}, Guillemin \cite{Guillemin-Weyl,Guillemin-residue}
and Wodzicki \cite{Wodzicki-residue,Wodzicki-cyclic}, we now look for traces on
the algebra $\Psiisoph{\Z}(\R^n)$. For $A\in\Psiiso{-\infty}$, by
\eqref{determinants:eq:trace-symbol}, $\Tr(A)=(2\pi)^{-n}\int a_L(x,\xi)\,dx\,d\xi$
is the integral of the left symbol; this is also true for any amplitude: if
$A=\Op(a)$ with $a(x,x',\xi)$ an amplitude of order $-\infty$, then
$\Tr(A)=(2\pi)^{-n}\int a(x,x,\xi)\,dx\,d\xi$. The integral extends to
$\Psiiso{m}$ for $m<-2n$, as we saw, but not beyond. To go further we
regularize it by embedding operators in holomorphic families, which moves us
from the symbol classes $\Gs{m}$ to the polyhomogeneous operators.

\begin{lemma}\label{determinants:lem:meromorphic}
  Let $\Omega\subset\C$ be open and connected, $m\in\C$, and $A(z)$ a
  holomorphic family of order $m+z$ on $\Omega$, with $b$, $b_j$ as in
  \eqref{determinants:eq:hol-family}. Suppose $\Omega$ contains points with
  $\Re z<-2n-\Re m$. Then $\Tr(A(z))$, defined where $\Re z<-2n-\Re m$, extends
  to a meromorphic function on $\Omega$ with at most simple poles, at the
  points $z_j=j-2n-m\in\Omega$, $j\in\N_0$, and
  \begin{equation}\label{determinants:eq:residues}
    \Res_{z=z_j}\Tr(A(z))=-(2\pi)^{-n}\int_{\Sphere^{2n-1}}\tr b_j(z_j;\theta)\,d\theta .
  \end{equation}
  If $A(z)=A_y(z)$ depends smoothly on a parameter $y$, so do the residues and
  the coefficients of the Laurent expansions.
\end{lemma}
\begin{proof}
  For $\Re z<-2n-\Re m$ we have $\Tr(A(z))=f_1(z)+f_2(z)$ with
  \begin{gather*}
    f_1(z)=(2\pi)^{-n}\int_{\abs{Z}\le1}\tr a(z;Z)\,dZ,\\
    f_2(z)=(2\pi)^{-n}\int_{\Sphere^{2n-1}}\int_1^\infty r^{m+z+2n-1}\tr b(z;1/r,\theta)\,dr\,d\theta .
  \end{gather*}
  The first is holomorphic on $\Omega$. Write
  $b(z;t,\theta)=\sum_{j<k}t^jb_j(z;\theta)+t^kb_{(k)}(z;t,\theta)$, with
  $b_{(k)}$ smooth and holomorphic in $z$ (Taylor's formula with integral
  remainder). Since $\int_1^\infty r^{m+z+2n-1-j}\,dr=(j-2n-m-z)^{-1}$ for
  $\Re z<j-2n-\Re m$,
  \begin{equation*}
    f_2(z)=(2\pi)^{-n}\sum_{j<k}\frac{1}{j-2n-m-z}\int_{\Sphere^{2n-1}}\tr b_j(z;\theta)\,d\theta+f_2^{(k)}(z),
  \end{equation*}
  where $f_2^{(k)}(z)=(2\pi)^{-n}\int_{\Sphere^{2n-1}}\int_1^\infty r^{m+z+2n-1-k}\tr b_{(k)}(z;1/r,\theta)\,dr\,d\theta$
  is holomorphic in $\Re z<k-2n-\Re m$, the integral converging locally
  uniformly there. This identity holds pointwise for $\Re z<-2n-\Re m$, and its
  right side is meromorphic on $\Omega_k=\{z\in\Omega;\Re z<k-2n-\Re m\}$; the
  right sides for $k$ and $k+1$ agree on $\Omega_k$ (the difference is the
  $j=k$ term, computed by the same integral). Since $\Omega$ is connected, these
  define the meromorphic continuation, with the stated poles and residues. All
  the expressions depend smoothly on parameters if $b$ does.
\end{proof}

We continue to write $\Tr(A(z))$ for the meromorphic continuation. By
uniqueness of continuation, any identity between such functions which holds
for $\Re z\ll0$ holds on $\Omega$. For a holomorphic family $A(z)$ of order
$m+z$ we have $b_j(0;\theta)$ equal to the coefficient of the term of
homogeneity $m-j$ in the symbol of $A(0)$. Taking $j=2n+m$, the residue at
$z=0$ only depends on $A(0)$, and it vanishes unless $m\in\Z$, $m\ge-2n$.

\begin{definition}\label{determinants:def:residue-trace}
  For $A\in\Psiisoph{m}(\R^n;\C^N)$, $m\in\C$, the \emph{residue trace} is
  \begin{equation*}
    \TrRes(A)=-\Res_{z=0}\Tr(A(z)),
  \end{equation*}
  where $A(z)$ is any holomorphic family of order $m+z$, on a connected
  neighbourhood of $0$ containing points with $\Re z<-2n-\Re m$, with
  $A(0)=A$; for instance $A(z)=AH(z)$, with $H(z)$ from
  \eqref{determinants:eq:standard-regularizer}.
\end{definition}

\begin{proposition}\label{determinants:prop:residue-trace}
  The residue trace is well defined, and
  \begin{equation}\label{determinants:eq:residue-formula}
    \TrRes(A)=(2\pi)^{-n}\int_{\Sphere^{2n-1}}\tr a_{-2n}(\theta)\,d\theta,
  \end{equation}
  where $a_{-2n}$ is the term of homogeneity $-2n$ in the expansion of the
  left symbol of $A$ (so $\TrRes$ vanishes unless $m\in\Z$, $m\ge-2n$; in
  particular on $\Psiisoph{-2n-1}$). It is a trace on $\Psiisoph{\Z}$:
  \begin{equation}\label{determinants:eq:residue-trace-property}
    \TrRes([A,B])=0\quad\forall\,A,B\in\Psiisoph{\Z}(\R^n;\C^N).
  \end{equation}
\end{proposition}
\begin{proof}
  The formula is \eqref{determinants:eq:residues} with $j=2n+m$, and shows that
  $\TrRes(A)$ is independent of the family. If $A\in\Psiisoph{m_1}$,
  $B\in\Psiisoph{m_2}$, then $C(z)=[AH(z),B]$ is an entire family of order
  $m_1+m_2+z$ with $C(0)=[A,B]$. For $\Re z<-2n-m_1-m_2$,
  $\Tr(C(z))=\Tr(AH(z)B)-\Tr(BAH(z))=0$ by
  \cref{determinants:lem:trace-property}, since $AH(z)\in\Psiiso{m_1+\Re z}$.
  Hence the continuation vanishes identically and $\TrRes([A,B])=0$.
\end{proof}

\begin{remark}\label{determinants:rem:residue-other-symbols}
  The left symbol in \eqref{determinants:eq:residue-formula} may be replaced
  by the right symbol (or the Weyl symbol). Indeed the right symbol of $A$ has
  the expansion $a_R\sim\sum_\alpha\frac1{\alpha!}(-D_\xi)^\alpha\pa_x^\alpha a_L$,
  and each term with $\alpha\ne0$ is a derivative $\pa_{Z_k}g$ of a
  polyhomogeneous function. If $g$ is homogeneous of degree $1-2n$ then,
  integrating $\pa_{Z_k}g$ over $1\le\abs{Z}\le R$ and using the divergence
  theorem,
  \begin{equation*}
    \log R\int_{\Sphere^{2n-1}}\pa_{Z_k}g\,d\theta
    =\int_{\Sphere^{2n-1}}\bigl(R^{1-2n}R^{2n-1}-1\bigr)g(\theta)\theta_k\,d\theta=0,
  \end{equation*}
  so the homogeneous term of degree $-2n$ of each such derivative integrates to
  zero over the sphere. The residue trace is, up to a constant multiple, the
  unique trace on $\Psiisoph{\Z}(\R^n)$; this is the isotropic analogue of
  Wodzicki's theorem \cite{Wodzicki-residue}, and we shall neither prove nor
  use it.
\end{remark}

\begin{discussion}
Why does the residue trace see only one term of the symbol? Integrating a
symbol term of homogeneity $\mu$ over $|Z|\ge1$ in $\R^{2n}$ gives
$\int_1^\infty r^{\mu+2n-1}\,dr$. This converges for $\mu<-2n$, diverges
like a power for $\mu>-2n$, and diverges \emph{logarithmically} exactly when
$\mu=-2n$. Inserting the regularizer $\ang{Z}^z$ turns a power divergence
into a pole at a non-zero value of $z$, harmless at $z=0$, while the
logarithmic divergence becomes a pole at $z=0$. So the residue at $z=0$
records the coefficient of the logarithmic divergence, which is the
integral of the homogeneous term of degree $-2n$ over the sphere.
\end{discussion}

\begin{question}
Every residual operator, indeed every operator of order less than $-2n$, has
residue trace zero. Is the residue trace then useless on trace class
operators?
\end{question}
\begin{answer}
Yes, and this is the point. On operators of very negative order the
ordinary trace is available, and the residue trace vanishes because
$\Tr(A(z))$ has no pole at $z=0$. The residue trace is a trace on the
\emph{quotient} $\Psiisoph{\Z}/\Psiiso{-\infty}$, where no non-zero
multiple of the ordinary trace makes sense. It complements the ordinary
trace rather than extending it.
\end{answer}

\section{Exterior derivation}\label{determinants:sec:derivation}

The regularized trace of the next section depends on a choice of
regularizing family $A(z)$, and its failure to be a trace is governed by the
derivative $\dot A$ of the family at $z=0$. Commutation with $\dot A$, which
is essentially the logarithm of an elliptic operator of order one, is the
``exterior derivation'' introduced here.

\begin{definition}\label{determinants:def:regularizer}
  A \emph{regularizing family} is an entire holomorphic family $A(z)$ of order
  $z$, with values in $\Psiisoph{z}(\R^n;\C^N)$, with scalar principal symbol and
  $A(0)=\Id$. Its derivative at $0$ is denoted $\dot A=\frac{d}{dz}A(z)|_{z=0}$.
\end{definition}

If $\sigma_L(A(z))=\ang{Z}^z\tilde a(z;Z)$ with $\tilde a$ holomorphic in $z$
with values in $\CI(\RC{\R^{2n}})$, then $A(0)=\Id$ means $\tilde a(0)=1$, so
\begin{equation*}
  \dot A=\Op\bigl(\log\ang{Z}+c\bigr),\qquad c=\pa_z\tilde a(0)\in\Symph{0}(\R^{2n})\text{ with scalar leading term.}
\end{equation*}
Thus $\dot A\in\Psiiso{\eps}$ for every $\eps>0$, but $\dot A\notin\Psiisoph{\Z}$
because of the logarithm. For the standard family $H(z)$ of
\eqref{determinants:eq:standard-regularizer}, $\dot H=\Op(\log\ang{Z})=\Op(\tfrac12\log(1+\abs{x}^2+\abs{\xi}^2))$.

\begin{proposition}\label{determinants:prop:derivation}
  For a regularizing family $A(z)$, the map
  \begin{equation*}
    D_AB=[\dot A,B]
  \end{equation*}
  is a derivation of $\Psiisoph{\Z}(\R^n;\C^N)$ with
  $D_A:\Psiisoph{m}\to\Psiisoph{m-1}$ and $D_A(\Psiiso{-\infty})\subset\Psiiso{-\infty}$.
  Equivalently, $D_AB=\frac{d}{dz}\bigl(A(z)BA(z)^{-1}\bigr)\big|_{z=0}$. For two
  regularizing families, $D_{A_2}-D_{A_1}=[\dot A_2-\dot A_1,\cdot\,]$ is an
  inner derivation, given by commutation with an element of $\Psiisoph{0}$.
  Moreover
  \begin{equation}\label{determinants:eq:res-derivation}
    \TrRes(D_AB)=0\quad\forall\,B\in\Psiisoph{\Z}(\R^n;\C^N).
  \end{equation}
\end{proposition}
\begin{proof}
  A commutator is a derivation. For $B\in\Psiisoph{m}$, $K(z)=[A(z),B]/z$ is an
  entire holomorphic family of order $m+z$ (the numerator is a holomorphic
  family vanishing at $z=0$, and dividing its normalized symbol by $z$
  preserves holomorphy), with $K(0)=[\dot A,B]$; hence $D_AB\in\Psiisoph{m}$.
  Since the principal symbol of $A(z)$ is scalar, $[A(z),B]$ has order
  $m+z-1$, so in fact $D_AB\in\Psiisoph{m-1}$. Residual operators form an ideal
  in $\bigcup_m\Psiiso{m}$, so $D_A$ preserves them. Near $z=0$, $A(z)$ is invertible
  with holomorphic inverse by \cref{determinants:lem:local-inversion}, and
  $\frac{d}{dz}(A(z)BA(z)^{-1})|_{z=0}=\dot AB-B\dot A$ by the product rule.
  The statement about two families follows since
  $\dot A_2-\dot A_1=\Op(c_2-c_1)\in\Psiisoph{0}$. Finally, $K(z)$ has value
  $D_AB$ at $0$ and $\Tr(K(z))=z^{-1}\bigl(\Tr(A(z)B)-\Tr(BA(z))\bigr)=0$ for
  $\Re z\ll0$ by \cref{determinants:lem:trace-property}, so the residue at
  $0$ vanishes.
\end{proof}

\begin{remark}\label{determinants:rem:exterior}
  $D_H$ is not an inner derivation of $\Psiisoph{\Z}$, which is why we call it
  an \emph{exterior} derivation. Suppose $[\dot H-P,B]=0$ for all
  $B\in\Psiisoph{\Z}$, for some $P\in\Psiisoph{\Z}$, and let $c$ be the left symbol
  of $C=\dot H-P$. The composition formula is exact when one factor has
  polynomial symbol, and gives $[C,x_j]=\Op(-i\pa_{\xi_j}c)$ and
  $[C,D_j]=\Op(-D_{x_j}c)$; hence $c$ is constant, which is impossible since
  $c-\log\ang{Z}$ is polyhomogeneous while $\log\ang{Z}$ is not.
\end{remark}

\section{The regularized trace}\label{determinants:sec:regularized}

Let $A(z)$ be a regularizing family. For $B\in\Psiisoph{\Z}(\R^n;\C^N)$ the
function $\Tr(BA(z))$ is meromorphic on $\C$ with at most a simple pole at $0$,
with residue $-\TrRes(B)$. We define the \emph{regularized trace} as the
finite part at $z=0$, the constant term of the Laurent expansion:
\begin{equation}\label{determinants:eq:regularized-trace}
  \TrReg{A}(B)=\FP_{z=0}\Tr(BA(z))=\lim_{z\to0}\Bigl(\Tr(BA(z))+\frac{\TrRes(B)}{z}\Bigr).
\end{equation}
If $B$ has order less than $-2n$ (in particular if $B$ is residual), then
$\Tr(BA(z))$ is holomorphic near $0$ and $\TrReg{A}(B)=\Tr(B)$; thus $\TrReg{A}$
extends the trace. By \cref{determinants:lem:meromorphic}, if $B_y$ is a
smooth family then $\TrReg{A}(B_y)$ is smooth in $y$ and
$d\,\TrReg{A}(B_y)=\TrReg{A}(dB_y)$. In contrast with the residue trace,
$\TrReg{A}$ depends on the choice of regularizing family, though only through
$\dot A$.

\begin{lemma}\label{determinants:lem:regularizer-change}
  If $A_1(z)$, $A_2(z)$ are regularizing families, then for
  $B\in\Psiisoph{\Z}(\R^n;\C^N)$,
  \begin{equation*}
    \TrReg{A_2}(B)-\TrReg{A_1}(B)=-\TrRes\bigl(B(\dot A_2-\dot A_1)\bigr).
  \end{equation*}
\end{lemma}
\begin{proof}
  $L(z)=(A_2(z)-A_1(z))/z$ is an entire family of order $z$ with
  $L(0)=\dot A_2-\dot A_1\in\Psiisoph{0}$, and
  $\Tr(BA_2(z))-\Tr(BA_1(z))=z\Tr(BL(z))$. The function $\Tr(BL(z))$ has at
  most a simple pole at $0$ with residue $-\TrRes(BL(0))$, which is therefore
  the value at $0$ of $z\Tr(BL(z))$.
\end{proof}

In particular, if $B$ has non-integral order, $\TrReg{A}(B)$ is independent of
$A$. The regularized trace is \emph{not} a trace; its failure to be one is
measured by the residue trace and the exterior derivation.

\begin{discussion}
The trace defect formula below is easy to remember. For $\Re z\ll0$ the
trace property gives $\Tr([B_1,B_2]A(z))=\Tr(B_1[B_2,A(z)])$. The
commutator $[B_2,A(z)]$ vanishes at $z=0$, so it is $z$ times something,
and that something is $-D_AB_2$ at $z=0$. A factor $z$ times a function with
a simple pole leaves a finite value, the residue. So the defect is a
residue trace: it is computed from finitely many terms of the symbols and
is local.
\end{discussion}

\begin{proposition}[Trace defect formula]\label{determinants:prop:trace-defect}
  For a regularizing family $A(z)$ and $B_1,B_2\in\Psiisoph{\Z}(\R^n;\C^N)$,
  \begin{equation}\label{determinants:eq:trace-defect}
    \TrReg{A}([B_1,B_2])=\TrRes(B_1D_AB_2)=-\TrRes(B_2D_AB_1).
  \end{equation}
\end{proposition}
\begin{proof}
  For $\Re z\ll0$, \cref{determinants:lem:trace-property} gives
  $\Tr(B_2B_1A(z))=\Tr(B_1A(z)B_2)$, so
  \begin{equation*}
    \Tr([B_1,B_2]A(z))=\Tr\bigl(B_1[B_2,A(z)]\bigr)=-z\Tr(B_1K(z)),
  \end{equation*}
  where $K(z)=[A(z),B_2]/z$ is the entire family of order $m_2+z$ with
  $K(0)=D_AB_2$ used in the proof of \cref{determinants:prop:derivation}. By
  continuation this holds for all $z$; since $\Tr(B_1K(z))$ has at most a
  simple pole at $0$ with residue $-\TrRes(B_1D_AB_2)$, the right side is
  regular at $0$ with value $\TrRes(B_1D_AB_2)$. The second equality follows
  from \eqref{determinants:eq:res-derivation} and the trace property:
  $0=\TrRes(D_A(B_1B_2))=\TrRes(B_2D_AB_1)+\TrRes(B_1D_AB_2)$.
\end{proof}

\begin{corollary}\label{determinants:cor:reg-conjugation}
  If $G\in\Giso{-\infty}(\R^n;\C^N)$ and $B\in\Psiisoph{\Z}(\R^n;\C^N)$, then
  $\TrReg{A}(G^{-1}BG)=\TrReg{A}(B)$.
\end{corollary}
\begin{proof}
  $\TrReg{A}(G^{-1}BG)-\TrReg{A}(BGG^{-1})=\TrReg{A}([G^{-1},BG])=-\TrRes(BG\,D_A(G^{-1}))$,
  which vanishes since $D_A(G^{-1})=D_A(G^{-1}-\Id)$ is residual and $\TrRes$
  vanishes on residual operators.
\end{proof}

\section{The index formula}\label{determinants:sec:index-formula}

Let $D_H$ be the exterior derivation of the standard family
\eqref{determinants:eq:standard-regularizer}. It preserves $\Psiiso{-\infty}$,
and $\TrRes$ vanishes there, so both descend to the quotient algebra
$\mathcal A^{\Z}=\Psiisoph{\Z}/\Psiiso{-\infty}$:
$D_H:\mathcal A^m\to\mathcal A^{m-1}$ and $\TrRes:\mathcal A^{\Z}\to\C$.

\begin{theorem}\label{determinants:thm:index-formula}
  If $A\in\Psiisoph{m}(\R^n;\C^N)$, $m\in\Z$, is elliptic with parametrix $B$, then
  \begin{equation}\label{determinants:eq:index-residue}
    \ind(A)=\TrRes(AD_HB)=-\TrRes(BD_HA).
  \end{equation}
  Thus, for invertible $a\in\mathcal A^{\Z}$, $\ind(a)=\TrRes(a\,D_H(a^{-1}))$.
  The index is locally constant along smooth families of elliptic operators,
  and $\ind(a_1a_2)=\ind(a_1)+\ind(a_2)$.
\end{theorem}
\begin{proof}
  Since $[A,B]$ is residual,
  $\ind(A)=\Tr([A,B])=\TrReg{H}([A,B])=\TrRes(AD_HB)$ by
  \eqref{determinants:eq:index-trace} and
  \eqref{determinants:eq:trace-defect}. If $A_t$ is a smooth family with smooth
  family of parametrices $B_t$ \ref{determinants:I:families}, the right side
  of \eqref{determinants:eq:index-residue} is an integral of finitely many
  terms of the symbol expansions, so it is continuous in $t$; being
  integer-valued it is locally constant. (Alternatively, $[A_t,B_t]$ is a
  smooth residual family.) For additivity,
  \begin{equation*}
    a_1a_2\,D_H(a_2^{-1}a_1^{-1})=a_1\bigl(a_2D_H(a_2^{-1})\bigr)a_1^{-1}+a_1D_H(a_1^{-1}),
  \end{equation*}
  and the residue trace of the first term is $\TrRes(a_2D_H(a_2^{-1}))$ by
  the trace property.
\end{proof}

Both properties hold for general Fredholm operators, but the point here is
that they follow from a simple algebraic formula.

\section{The variation $1$-form}\label{determinants:sec:variation}

The index is the obstruction to perturbing a single elliptic operator, by a
smoothing operator, to be invertible. We now derive a closely related
obstruction for families. Let $Y$ be a manifold and
\begin{equation}\label{determinants:eq:family}
  Y\ni y\longmapsto A_y\in\Ell^0_{\mathrm{iso},0}(\R^n;\C^N)
\end{equation}
a smooth family of elliptic operators of order $0$ and index $0$. The
\emph{families perturbative invertibility question} asks whether there is a
smooth family
\begin{equation}\label{determinants:eq:perturbation}
  Y\ni y\longmapsto E_y\in\Psiiso{-\infty}(\R^n;\C^N)\quad\text{with}\quad A_y+E_y\in\Giso{0}(\R^n;\C^N)\ \forall\,y.
\end{equation}
Since the index vanishes, such $E_y$ exists for each fixed $y$
(\cref{determinants:prop:invertible-perturbation}), and also locally: if
$A_{y_0}+F$ is invertible then so is $A_y+F$ for $y$ near $y_0$, with inverse
a smooth family in $\Psiisoph{0}$ by \cref{determinants:lem:local-inversion}.
The issue is the global existence of a smooth (or even continuous) choice.

\begin{question}
If every $A_y$ can be made invertible by a smoothing perturbation, why can
we not just patch the local choices together with a partition of unity?
\end{question}
\begin{answer}
Invertibility is not preserved by taking convex combinations: if $A_y+F_1$
and $A_y+F_2$ are invertible, $A_y+\frac12(F_1+F_2)$ need not be. The local
choices can be compared only through the group $\Giso{-\infty}$, by
$A_y+F_2=(A_y+F_1)(\Id+X_y)$, and the obstruction to patching them lives in
the topology of the transition functions $\Id+X_y$. Taking determinants of
these turns them into $\C^*$-valued functions, that is, into the transition
functions of a line bundle, and the obstruction becomes its first Chern
class (\cref{determinants:sec:det-bundle}).
\end{answer}

Operator-valued differential forms are handled coefficientwise: if
$\omega=\sum\omega_j\,dy_j$ and $\eta=\sum\eta_k\,dy_k$ in local coordinates,
$\omega\wedge\eta=\sum_{j,k}\omega_j\eta_k\,dy_j\wedge dy_k$, and $\TrRes$ and
$\TrReg{H}$ are applied coefficientwise. Let $a_y=[A_y]\in\mathcal A^0$ and let
$B_y$ be a smooth family of parametrices, so $a_y^{-1}=[B_y]$. We set
\begin{equation*}
  \theta=a^{-1}da,\qquad \varphi=a^{-1}D_Ha,
\end{equation*}
represented by $B_y\,dA_y$ and $B_yD_HA_y$, and define
\begin{equation}\label{determinants:eq:beta}
  \beta=\tfrac12\TrRes\bigl(a^{-1}da\wedge a^{-1}da\;a^{-1}D_Ha\bigr)=\tfrac12\TrRes(\theta\wedge\theta\,\varphi),\qquad
  \gamma=\tfrac12\TrRes(\theta\varphi).
\end{equation}
These are smooth forms on $Y$ of degrees $2$ and $1$; they depend only on
$a_y$, since $\TrRes$ vanishes on residual operators.

\begin{proposition}\label{determinants:prop:variation}
  The $2$-form $\beta$ is closed. If there is a smooth family $E_y$ as in
  \eqref{determinants:eq:perturbation}, then, with $\tilde A_y=A_y+E_y$ and
  \begin{equation}\label{determinants:eq:alpha-bar}
    \bar\alpha=\TrReg{H}\bigl(\tilde A_y^{-1}d\tilde A_y\bigr),
  \end{equation}
  we have $d\bar\alpha=-\beta+d\gamma$. In particular $\beta$ is exact, so its
  de Rham class $[\beta]\in H^2(Y;\C)$ is an obstruction to the existence of
  a smooth invertible perturbation.
\end{proposition}
\begin{proof}
  Suppose first that $\tilde A=A+E$ is invertible on $Y$, and compute with
  $\theta=\tilde A^{-1}d\tilde A$ and $\varphi=\tilde A^{-1}D_H\tilde A$,
  which represent the classes above. Since $d(\tilde A^{-1})=-\tilde A^{-1}d\tilde A\,\tilde A^{-1}$,
  \begin{equation*}
    d\theta=-\theta\wedge\theta=-\tfrac12\sum_{j,k}[\theta_j,\theta_k]\,dy_j\wedge dy_k .
  \end{equation*}
  Since $d$ commutes with $\TrReg{H}$, the trace defect formula
  \eqref{determinants:eq:trace-defect} gives
  \begin{equation*}
    d\bar\alpha=-\tfrac12\sum_{j,k}\TrRes(\theta_jD_H\theta_k)\,dy_j\wedge dy_k=-\tfrac12\TrRes(\theta\wedge D_H\theta).
  \end{equation*}
  As $D_H$ is a derivation commuting with $d$ (it is commutation with the fixed
  operator $\dot H$), $D_H(\tilde A^{-1})=-\tilde A^{-1}(D_H\tilde A)\tilde A^{-1}$
  and
  \begin{equation*}
    D_H\theta=-\varphi\theta+\tilde A^{-1}d(D_H\tilde A)=d\varphi+\theta\varphi-\varphi\theta .
  \end{equation*}
  By the trace property, $\TrRes(\theta\wedge\varphi\theta)=\sum\TrRes(\theta_k\theta_j\varphi)\,dy_j\wedge dy_k=-\TrRes(\theta\wedge\theta\varphi)$,
  so $\TrRes(\theta\wedge D_H\theta)=\TrRes(\theta\wedge d\varphi)+2\TrRes(\theta\wedge\theta\varphi)$.
  On the other hand $d(\theta\varphi)=d\theta\,\varphi-\theta\wedge d\varphi$, so
  $2\,d\gamma=-\TrRes(\theta\wedge\theta\varphi)-\TrRes(\theta\wedge d\varphi)$.
  Combining,
  \begin{equation*}
    \TrRes(\theta\wedge D_H\theta)=2\beta-2\,d\gamma,\qquad d\bar\alpha=-\beta+d\gamma .
  \end{equation*}
  Thus $\beta=d(\gamma-\bar\alpha)$ is exact. In general, invertible
  perturbations exist locally on $Y$, as noted above, so $\beta$ is locally
  exact and hence closed.
\end{proof}

The form $\gamma$ is defined on all of $Y$ for any family
\eqref{determinants:eq:family}, so $[\beta]=[\beta-d\gamma]$. On $\Giso{0}$,
$\bar\alpha$ extends the $1$-form $d\log\detFr=\Tr(A^{-1}dA)$ from
$\Giso{-\infty}$, but it is not closed in general, because the regularized
trace is not a trace. To explain the topological meaning of $[\beta]$ we now
construct the determinant line bundle.

\section{The determinant bundle}\label{determinants:sec:det-bundle}

The determinant bundle was introduced by Quillen \cite{Quillen} for families of
Cauchy--Riemann operators. Recall that $\detFr$ is a character of
$\Giso{-\infty}$: $\detFr(AB)=\detFr(A)\detFr(B)$. Instead of trying to extend
it multiplicatively to $\Giso{0}$ (see \cref{determinants:rem:no-extension}),
we consider all its possible values at once.

\begin{discussion}
An invertible $A\in G^0$ has no preferred determinant. But if we change $A$
to $A(\Id+E)$ with $\Id+E\in G^{-\infty}$, the ``determinant'' ought to be
multiplied by $\detFr(\Id+E)$. The relation
\eqref{determinants:eq:det-equivalence} below records exactly this rule.
Over each class $a=[A]$ the possible values then form a complex line, not a
number. As $a$ varies these lines fit together into a line bundle, and the
question whether a multiplicative determinant exists becomes the question
whether this bundle is trivial.
\end{discussion}

Write
$G^0=\Giso{0}(\R^n;\C^N)$ and $G^{-\infty}=\Giso{-\infty}(\R^n;\C^N)$; the
latter is a normal subgroup of $G^0$, and $\mathcal G=G^0/G^{-\infty}$ is the
set of classes $a=[A]\in\mathcal A^0$ of invertible operators, which by
\cref{determinants:prop:invertible-perturbation} are exactly the invertible
elements of $\mathcal A^0$ of index zero. On $G^0\times\C$ define
\begin{equation}\label{determinants:eq:det-equivalence}
  (A,z)\sim(A',z')\iff A'=A(\Id+E),\ z'=\detFr(\Id+E)\,z\ \text{ for some }\Id+E\in G^{-\infty}.
\end{equation}

\begin{lemma}\label{determinants:lem:det-line}
  \eqref{determinants:eq:det-equivalence} is an equivalence relation. Let
  $\mathcal D=(G^0\times\C)/{\sim}$ and $\pi:\mathcal D\to\mathcal G$,
  $\pi[(A,z)]=[A]$. Each fibre $\mathcal D_a=\pi^{-1}(a)$ is a one-dimensional
  complex vector space with the operations
  $[(A,z)]+[(A,w)]=[(A,z+w)]$, $\lambda[(A,z)]=[(A,\lambda z)]$.
\end{lemma}
\begin{proof}
  Reflexivity, symmetry and transitivity follow from $\detFr(\Id)=1$,
  $\detFr((\Id+E)^{-1})=\detFr(\Id+E)^{-1}$ and multiplicativity. If
  $[A]=[A_0]=a$ then $A_0^{-1}A\in G^0\cap(\Id+\Psiiso{-\infty})=G^{-\infty}$,
  so every element of $\mathcal D_a$ has a representative $(A_0,w)$, and $w$
  is unique since $(A_0,w)\sim(A_0,w')$ forces $\Id+E=\Id$. Thus
  $[(A_0,w)]\mapsto w$ is a bijection $\mathcal D_a\to\C$. The operations
  above are well defined because the relation is linear in $z$, and they
  correspond to those of $\C$ under this bijection.
\end{proof}

Now let $y\mapsto A_y$ be a smooth family \eqref{determinants:eq:family} over a
manifold $Y$ and $a_y=[A_y]$. Set $\mathcal D_Y=\bigsqcup_{y\in Y}\mathcal D_{a_y}$.
A \emph{local invertible lift} on $U\subset Y$ open is a smooth family
$\tilde A_y=A_y+E_y\in G^0$, $y\in U$, with $E_y$ residual; it defines the
section $s_{\tilde A}(y)=[(\tilde A_y,1)]$ of $\mathcal D_Y$ over $U$. Local
invertible lifts exist near every point of $Y$.

\begin{proposition}\label{determinants:prop:det-bundle}
  $\mathcal D_Y$ has a unique structure of smooth complex line bundle over $Y$
  for which the sections $s_{\tilde A}$ are smooth. If $\tilde A'=\tilde A(\Id+X)$
  is another local invertible lift, then $X_y$ is a smooth residual family and
  \begin{equation}\label{determinants:eq:transition}
    s_{\tilde A'}=\detFr(\Id+X)^{-1}s_{\tilde A}.
  \end{equation}
  If \eqref{determinants:eq:perturbation} has a solution, $\mathcal D_Y$ is
  trivial.
\end{proposition}
\begin{proof}
  $\Id+X=\tilde A^{-1}\tilde A'=\Id+\tilde A^{-1}(E'-E)$, and $\tilde A^{-1}$
  is a smooth family in $\Psiisoph{0}$ by
  \cref{determinants:lem:local-inversion}, so $X$ is a smooth residual family.
  By \eqref{determinants:eq:det-equivalence},
  $(\tilde A,\detFr(\Id+X)^{-1})\sim(\tilde A(\Id+X),1)$, which is
  \eqref{determinants:eq:transition}. The transition functions
  $\detFr(\Id+X_y)^{-1}$ are smooth and non-vanishing by
  \cref{determinants:thm:det-entire}, and the cocycle condition holds
  automatically since all sections live in the same space $\mathcal D_Y$.
  Hence the local trivializations $s_{\tilde A}$ define a smooth line bundle
  structure, which is unique since a line bundle structure is determined by
  a family of local non-vanishing smooth sections. A global invertible lift gives a global
  non-vanishing section.
\end{proof}

The projection to $\mathcal G$ and the quotient map fit into the commutative
diagram
\begin{equation*}
  \begin{tikzcd}
    G^0\times\C\arrow[r,"/\sim"]\arrow[d] & \mathcal D\arrow[d,"\pi"]\\
    G^0\arrow[r,"/G^{-\infty}"] & \mathcal G,
  \end{tikzcd}
\end{equation*}
and the pull-back of $\mathcal D$ to $G^0$ is canonically trivial, with the
section $A\mapsto[(A,1)]$. Over $Y$, the form $\bar\alpha$ of
\eqref{determinants:eq:alpha-bar} defines a connection on $\mathcal D_Y$.

\begin{theorem}\label{determinants:thm:connection}
  There is a unique connection $\nabla$ on $\mathcal D_Y$ such that, for every
  local invertible lift $\tilde A$,
  \begin{equation*}
    \nabla s_{\tilde A}=-\TrReg{H}(\tilde A^{-1}d\tilde A)\otimes s_{\tilde A}.
  \end{equation*}
  Its curvature is the globally defined $2$-form $\beta-d\gamma$ of
  \eqref{determinants:eq:beta}.
\end{theorem}
\begin{proof}
  Write $\bar\alpha_{\tilde A}=\TrReg{H}(\tilde A^{-1}d\tilde A)$. We must check
  that the local definitions agree. With $\tilde A'=\tilde A(\Id+X)$,
  \begin{equation*}
    \tilde A'^{-1}d\tilde A'=(\Id+X)^{-1}\bigl(\tilde A^{-1}d\tilde A\bigr)(\Id+X)+(\Id+X)^{-1}dX,
  \end{equation*}
  so by \cref{determinants:cor:reg-conjugation}, and since $\TrReg{H}=\Tr$ on
  residual operators,
  \begin{equation*}
    \bar\alpha_{\tilde A'}=\bar\alpha_{\tilde A}+\Tr((\Id+X)^{-1}dX)=\bar\alpha_{\tilde A}+d\log\detFr(\Id+X).
  \end{equation*}
  With $f=\detFr(\Id+X)^{-1}$, the Leibniz rule and
  \eqref{determinants:eq:transition} give
  \begin{equation*}
    \nabla s_{\tilde A'}=\nabla(fs_{\tilde A})=(d\log f-\bar\alpha_{\tilde A})\,s_{\tilde A'}=-\bar\alpha_{\tilde A'}\,s_{\tilde A'},
  \end{equation*}
  as required. The curvature is $-d\bar\alpha_{\tilde A}=\beta-d\gamma$ by
  \cref{determinants:prop:variation}.
\end{proof}

Thus $\beta-d\gamma$ is a curvature form, and by Chern--Weil theory the
class $\frac{i}{2\pi}[\beta]$ represents the first Chern class of
$\mathcal D_Y$ in $H^2(Y;\C)$; in particular it is the image of an integral
class. This is the topological origin of the obstruction in
\cref{determinants:prop:variation}: the family $A_y$ admits a smooth invertible
perturbation only if $\mathcal D_Y$ is trivial.

\begin{remark}\label{determinants:rem:no-extension}
  If $f:G^0\to\C^*$ were a multiplicative extension of $\detFr$ which is smooth
  along smooth families, then the map $[(A,z)]\mapsto z/f(A)$ would be well
  defined on $\mathcal D$, since
  \begin{equation*}
    \detFr(\Id+E)\,z/f(A(\Id+E))=z/f(A),
  \end{equation*}
  and would give a smooth non-vanishing
  linear functional on each $\mathcal D_Y$; hence every $\mathcal D_Y$ would be
  trivial and every $[\beta]$ would vanish. So any family with $[\beta]\ne0$
  shows that no such extension exists.
\end{remark}

\section{Exercises}\label{determinants:sec:exercises}

These exercises cover both this chapter and \cref{ch:isotropic}.

\begin{exercise}[Isotropic Sobolev spaces]\label{determinants:ex:sobolev}
  For $k\in\N_0$ define
  \begin{gather*}
    \Hs{k}_{\mathrm{iso}}(\R^n)=\{u\in L^2(\R^n);\ x^\alpha D_x^\beta u\in L^2(\R^n)\ \forall\,\abs{\alpha}+\abs{\beta}\le k\},\\
    \Hs{-k}_{\mathrm{iso}}(\R^n)=\Bigl\{u\in\cSp(\R^n);\ u=\sum_{\abs{\alpha}+\abs{\beta}\le k}x^\alpha D_x^\beta u_{\alpha\beta},\ u_{\alpha\beta}\in L^2(\R^n)\Bigr\}.
  \end{gather*}
  Show that if $A\in\Psiisoph{p}(\R^n)$, $p\in\Z$, then
  $A:\Hs{k}_{\mathrm{iso}}\to\Hs{k-p}_{\mathrm{iso}}$ for all $k\in\Z$. Deduce,
  using elliptic isotropic operators, that the definition
  \begin{equation*}
    \Hs{s}_{\mathrm{iso}}(\R^n)=\{u\in\cSp(\R^n);\ Au\in L^2(\R^n)\ \forall\,A\in\Psiisoph{-s}(\R^n)\},\quad s\in\R,
  \end{equation*}
  is consistent with the previous one, and that
  \begin{gather*}
    A\in\Psiisoph{M}(\R^n)\implies A:\Hs{s}_{\mathrm{iso}}\to\Hs{s-M}_{\mathrm{iso}},\qquad
    \bigcap_s\Hs{s}_{\mathrm{iso}}=\cS(\R^n),\quad\bigcup_s\Hs{s}_{\mathrm{iso}}=\cSp(\R^n),\\
    A\in\Psiisoph{m}(\R^n)\text{ elliptic},\ u\in\cSp(\R^n),\ Au\in\Hs{s}_{\mathrm{iso}}\implies u\in\Hs{s+m}_{\mathrm{iso}}.
  \end{gather*}
\end{exercise}

\begin{exercise}\label{determinants:ex:compact-inclusion}
  Show that if $\eps>0$ then
  $\Hs{\eps}_{\mathrm{iso}}(\R^n)\subset\ang{x}^{-\eps}L^2(\R^n)\cap\Hs{\eps}(\R^n)$,
  and deduce that the inclusion $\Hs{\eps}_{\mathrm{iso}}(\R^n)\hookrightarrow L^2(\R^n)$
  is compact.
\end{exercise}

\begin{exercise}\label{determinants:ex:compact}
  Using \cref{determinants:ex:compact-inclusion}, or otherwise, show that each
  element of $\Psiisoph{-\eps}(\R^n)$, $\eps>0$, is compact on $L^2(\R^n)$.
\end{exercise}

\begin{exercise}\label{determinants:ex:finite-rank-inverse}
  Without using the Fredholm theory of \ref{determinants:I:residual}, show that
  if $E\in\Psiiso{-\infty}(\R^n)$ there is $F\in\Psiiso{-\infty}(\R^n)$ such
  that $(\Id+E)(\Id+F)=\Id+G$ with $G\in\Psiiso{-\infty}(\R^n)$ of finite rank,
  i.e.\ $G\cdot\cS(\R^n)$ is finite dimensional. (Hint: write $E=E_1+E_2$ with
  $E_1$ of finite rank and $\norm{E_2}_{\Lin(L^2)}<1$, and use
  \cref{determinants:lem:residual-inverse}.)
\end{exercise}

\begin{exercise}\label{determinants:ex:fredholm}
  Using \cref{determinants:ex:finite-rank-inverse}, show that an elliptic
  $A\in\Psiisoph{m}(\R^n)$ has a parametrix $B\in\Psiisoph{-m}(\R^n)$ up to finite
  rank error, i.e.\ $AB-\Id$ and $BA-\Id$ are finite rank elements of
  $\Psiiso{-\infty}(\R^n)$. Deduce that $A$ is Fredholm as an operator
  $\Hs{M}_{\mathrm{iso}}(\R^n)\to\Hs{M-m}_{\mathrm{iso}}(\R^n)$ for every $M$.
  (An operator between Hilbert spaces is Fredholm if it is bounded, with finite
  dimensional null space and closed range of finite codimension.)
\end{exercise}

\begin{exercise}[The harmonic oscillator]\label{determinants:ex:oscillator}
  Show that $H=\abs{D}^2+\abs{x}^2=\sum_j(D_j^2+x_j^2)$ is an elliptic element of
  $\Psiisoph{2}(\R^n)$. With the creation and annihilation operators
  $C_j=D_j+ix_j$ and $A_j=D_j-ix_j=C_j^*$, show that
  \begin{gather*}
    H=\sum_{j=1}^nC_jA_j+n=\sum_{j=1}^nA_jC_j-n,\\
    [A_j,H]=2A_j,\quad[C_j,H]=-2C_j,\quad[C_l,C_j]=0,\quad[A_l,A_j]=0,\quad[A_l,C_j]=2\delta_{lj}\Id .
  \end{gather*}
  Knowing that $(H-\lambda)u=0$ with $u\in\cSp(\R^n)$ implies $u\in\cS(\R^n)$
  (why?), show that
  \begin{gather*}
    E_\lambda=\{u\in\cSp(\R^n);\ (H-\lambda)u=0\}\ne\{0\}\iff\lambda\in n+2\N_0,\\
    E_{n+2k}=\Bigl\{\sum_{\abs{\alpha}=k}c_\alpha C^\alpha e^{-\abs{x}^2/2};\ c_\alpha\in\C\Bigr\},\quad k\in\N_0 .
  \end{gather*}
\end{exercise}

\begin{exercise}[Determinants of matrices]\label{determinants:ex:matrix-det}
  Show that there is a unique function $\det:M(N;\C)\to\C$ which is linear in
  each column, vanishes when two columns are equal, and satisfies
  $\det\Id=1$, and that
  $\det A=\sum_{\sigma}\sgn(\sigma)\prod_iA_{i\sigma(i)}$. Show that
  $\det(AB)=\det A\det B$ (consider $A\mapsto\det(AB)$ for fixed $B$), that $A$
  is invertible if and only if $\det A\ne0$, that
  $\frac{d}{ds}\det(\Id+sA)|_{s=0}=\tr A$, and deduce Jacobi's formula
  \eqref{determinants:eq:jacobi}. Show also that $\det$ is the unique
  multiplicative polynomial function on $M(N;\C)$ with
  $\frac{d}{ds}\det(\Id+sA)|_{s=0}=\tr A$.
\end{exercise}

\begin{exercise}[The $1$-form $\alpha$ is closed]\label{determinants:ex:closed}
  Let $A(s,t)$ be a smooth two-parameter family in $\Giso{-\infty}(\R^n)$ and
  $\alpha=\Tr(A^{-1}\pa_sA)\,ds+\Tr(A^{-1}\pa_tA)\,dt$ the pull-back of
  $\Tr(A^{-1}dA)$. Show that
  \begin{equation*}
    d\alpha=\Tr\bigl([A^{-1}\pa_tA,A^{-1}\pa_sA]\bigr)\,ds\wedge dt,
  \end{equation*}
  and hence that $d\alpha=0$. Explain why this implies that
  $\Tr(A^{-1}dA)$ is closed on $\Giso{-\infty}$, and show directly, using
  \cref{determinants:lem:path-limit}, that its integral over any closed path
  lies in $2\pi i\Z$.
\end{exercise}

\begin{exercise}[Bundle coefficients]\label{determinants:ex:bundles}
  If $V$ and $W$ are finite dimensional complex vector spaces, show that the
  space $\Psiisoph{m}(\R^n;V,W)$ of operators from $\cS(\R^n;V)$ to
  $\cSp(\R^n;W)$ given by matrices with entries in $\Psiisoph{m}(\R^n)$, with
  respect to bases of $V$ and $W$, is independent of the choice of bases.
  Extend the results of \cref{ch:isotropic} to these spaces.
\end{exercise}

\begin{exercise}[Ellipticity is necessary for the Fredholm property]\label{determinants:ex:necessity}
  Let $A\in\Psiisoph{m}(\R^n;\C^N)$ be Fredholm as an operator
  $\Hs{s}_{\mathrm{iso}}\to\Hs{s-m}_{\mathrm{iso}}$ for some $s$. Show that $A$
  is elliptic, as follows.
  \begin{enumerate}
  \item Reduce to the case $m=s=0$, using invertible elliptic operators of
    appropriate orders.
  \item If the principal symbol of $A$ is not invertible at
    $\theta\in\Sphere^{2n-1}$, construct a sequence $u_k\in\cS$ with
    $\norm{u_k}_{L^2}=1$, $u_k\to0$ weakly and $Au_k\to0$ in $L^2$ (consider
    $u_k=e^{ik^2x\cdot\xi_0}\phi(x/k-kx_0)v$, suitably normalized, with
    $(x_0,\xi_0)=\theta$ and $v$ in the null space of the principal symbol at
    $\theta$).
  \item Show that a Fredholm operator satisfies $\norm{u}\le C\norm{Au}+C\norm{\Pi u}$ with $\Pi$ the
    projection onto its finite dimensional null space, and derive a
    contradiction.
  \end{enumerate}
\end{exercise}

\begin{exercise}[The Koszul complex]\label{determinants:ex:koszul}
  Let $\Lambda^k=\Lambda^k\C^n$ with basis $dx^\alpha=dx^{\alpha_1}\wedge\dots\wedge dx^{\alpha_k}$,
  $\alpha_1<\dots<\alpha_k$, declared orthonormal, and let
  $\cSp(\R^n;\Lambda^k)$ be the space of $u=\sum_\alpha u_\alpha\,dx^\alpha$,
  $u_\alpha\in\cSp(\R^n)$. Using the annihilation operators of
  \cref{determinants:ex:oscillator} define
  \begin{equation*}
    D:\cSp(\R^n;\Lambda^k)\to\cSp(\R^n;\Lambda^{k+1}),\qquad Du=\sum_{j=1}^n\sum_\alpha A_ju_\alpha\,dx^j\wedge dx^\alpha .
  \end{equation*}
  Show that $D^2=0$, and that the formal adjoint of $D$ with respect to the
  $L^2$ pairing is $D^*u=\sum_{j,\alpha}C_ju_\alpha\,\iota_j dx^\alpha$, where
  $\iota_j$ is contraction with $e_j$, the adjoint of $dx^j\wedge$. Show that
  $D+D^*\in\Psiisoph{1}(\R^n;\Lambda^*)$ is elliptic and that
  \begin{equation*}
    (D+D^*)^2=DD^*+D^*D=H-n+2\mathcal N,
  \end{equation*}
  where $\mathcal N=k$ on $\Lambda^k$. Deduce that the null space of $D+D^*$
  on $\cSp(\R^n;\Lambda^*)$ is spanned by $e^{-\abs{x}^2/2}\in\cS(\R^n;\Lambda^0)$,
  and that
  \begin{gather*}
    \{u\in\cSp(\R^n);\ Du=0\}=\C e^{-\abs{x}^2/2},\\
    \{u\in\cSp(\R^n;\Lambda^k);\ Du=0\}=D\,\cSp(\R^n;\Lambda^{k-1}),\ k\ge1 .
  \end{gather*}
  Conclude that $D+D^*$, as an elliptic element of
  $\Psiisoph{1}(\R^n;\Lambda^{\mathrm{even}},\Lambda^{\mathrm{odd}})$, has index
  $1$ (and as an operator from odd to even forms, index $-1$).
\end{exercise}

\begin{exercise}[Isotropic essential support]\label{determinants:ex:ess-supp}
  For $A\in\Psiisoph{m}(\R^n)$ with left symbol $a$, define the isotropic
  essential support (operator wavefront set) $\WFiso'(A)\subset\R^{2n}\setminus\{0\}$
  as the complement of the set of $Z_0\ne0$ having a conic neighbourhood
  $\Gamma$ in which $a$ is of order $-\infty$:
  $\abs{D^\alpha a(Z)}\le C_{\alpha,k}\ang{Z}^{-k}$ for $Z\in\Gamma$ and all
  $\alpha,k$. Show that $\WFiso'(A)$ is closed and conic, that it is unchanged if
  the left symbol is replaced by the right or Weyl symbol, and that
  \begin{gather*}
    \WFiso'(A+B)\subset\WFiso'(A)\cup\WFiso'(B),\qquad\WFiso'(AB)\subset\WFiso'(A)\cap\WFiso'(B),\\
    \WFiso'(A)=\emptyset\iff A\in\Psiiso{-\infty}(\R^n).
  \end{gather*}
\end{exercise}

\begin{exercise}[Isotropic partition of unity]\label{determinants:ex:partition}
  Let $U_i$ be a finite open cover of $\Sphere^{2n-1}$ and
  $\tilde U_i=\{Z\in\R^{2n}\setminus\{0\};\ Z/\abs{Z}\in U_i\}$ the
  corresponding conic cover of $\R^{2n}\setminus\{0\}$. Show that there are
  $A_i\in\Psiisoph{0}(\R^n)$ with $\WFiso'(A_i)\subset\tilde U_i$ and
  $\Id-\sum_iA_i\in\Psiiso{-\infty}(\R^n)$.
\end{exercise}

\begin{exercise}\label{determinants:ex:perturbation-sobolev}
  Suppose $A\in\Psiisoph{m}(\R^n)$ is elliptic and has index zero. Show that the
  operator $A+F$ of \cref{determinants:prop:invertible-perturbation} is an
  isomorphism of $\cS(\R^n)$, of $\cSp(\R^n)$, and from
  $\Hs{s}_{\mathrm{iso}}(\R^n)$ to $\Hs{s-m}_{\mathrm{iso}}(\R^n)$ for every
  $s$.
\end{exercise}

\begin{exercise}[Isotropic wavefront set]\label{determinants:ex:wf}
  For $A\in\Psiisoph{0}(\R^n)$ let
  $\Char_{\mathrm{iso}}(A)=\{Z\ne0;\ \sigma_0(A)(Z/\abs{Z})=0\}$, where
  $\sigma_0(A)$ is the principal symbol. For $u\in\cSp(\R^n)$ define
  \begin{equation*}
    \WFiso(u)=\bigcap\bigl\{\Char_{\mathrm{iso}}(A);\ A\in\Psiisoph{0}(\R^n),\ Au\in\cS(\R^n)\bigr\}.
  \end{equation*}
  Show that $\WFiso(u)$ is a closed cone, that $\WFiso(u)=\emptyset$ if and only
  if $u\in\cS(\R^n)$, and that $\WFiso(Bu)\subset\WFiso(u)\cap\WFiso'(B)$ for
  $B\in\Psiisoph{m}(\R^n)$. (Use \cref{determinants:ex:ess-supp,determinants:ex:partition}.)
\end{exercise}

\begin{exercise}[Traces and determinants of lower regularity]\label{determinants:ex:extend}
  Let $m<-2n$. Show that $\Tr$ is continuous on $\Psiiso{m}(\R^n)$, that
  $\Id+E$ with $E\in\Psiiso{m}(\R^n)$ invertible on $L^2$ has inverse
  $\Id+F$ with $F\in\Psiiso{m}(\R^n)$, and that the definition of $\detFr$
  by integration of $\Tr(A^{-1}dA)$ along paths extends to
  $\Id+\Psiiso{m}(\R^n)$ with the properties of
  \cref{determinants:prop:det-basic} and \cref{determinants:thm:det-entire}.
  (Hint: approximate $E$ by $\Op(\chi(Z/k)e)$ as in the proof of
  \cref{determinants:lem:trace-property}, which converges in
  $\Psiiso{m+\eps}$, and show that the corresponding determinants converge.)
\end{exercise}

\section*{Notes}

The main departures from the original notes are as follows. The Fredholm
determinant is constructed by integrating $\Tr(A^{-1}dA)$ along paths and
comparing with finite dimensional determinants via the harmonic oscillator
basis (\cref{determinants:lem:path-limit}); this gives path independence and
the limit formula \eqref{determinants:eq:det-limit} at once, and the
holomorphy across non-invertible operators is reduced, via
\cref{determinants:lem:reduction}, to a finite dimensional determinant. The
residue trace is normalized as $-\Res_{z=0}$, so that
\eqref{determinants:eq:residue-formula} holds with a positive sign; this agrees
with the usual convention for Wodzicki's residue \cite{Wodzicki-residue},
since $z$ here corresponds to $-s$ in the zeta-function regularization
$\Tr(AQ^{-s})$. The exterior derivation is defined directly as commutation
with $\dot A$, which avoids inverting the regularizing family away from $z=0$.
The determinant line bundle is constructed over finite dimensional families,
where it is an honest smooth line bundle, and we have added the connection of
\cref{determinants:thm:connection}, whose curvature is the obstruction form of
\cref{determinants:prop:variation}.

For the residue trace on compact manifolds see
\cite{Wodzicki-residue,Guillemin-residue}; for the isotropic (Shubin) calculus
and its spectral theory see \cite{Shubin}; the determinant bundle originates in
\cite{Quillen}. The Chern--Weil description of the first Chern class used at
the end of \cref{determinants:sec:det-bundle} can be found, for instance, in
Appendix~C of J.~Milnor and J.~Stasheff, \emph{Characteristic classes}.
